\section{The Fast Causal Inference Algorithm}\label{sec:fci}

\Joris{We should check what background knowledge we are now using for the input nodes. We're using JCI 1, but I'm not sure whether we're using JCI 2 as well. If we are, we should investigate what happens if we drop JCI 2.}

In this final chapter, we present an extension of the Fast Causal Inference (FCI) algorithm. 
The FCI algorithm is one of the highlights in the field of constraint-based causal discovery.
It was originally designed for purely observational (non-experimental) data and relied on the assumption of acyclicity, but allowed for latent variables (either marginalized over, or conditioned upon in case of selection bias) \cite{SMR1995,SMR1999}.
Later, the algorithm was augmented with additional `orientation rules' and this augmented FCI algorithm was shown to be complete in a certain sense \cite{Zhang2008_AI}.\footnote{The augmented version of \cite{Zhang2008_AI} is often referred to simply as `the FCI algorithm', and we will do so here as well.}
While the FCI algorithm was originally designed for the acyclic setting, it was recently discovered that it also works in case cycles are present (more specifically, for $\sigma$-faithful simple SCMs) \cite{MooijClaassen_UAI_20}.
While that work made a simplification by assuming no selection bias, we prove here that the FCI algorithm is sound even when the data is generated according to a conditional Markov kernel induced by a $\sigma$-faithful simple SCM, or in other words, that it can cope with the possible presence of cycles \emph{and} selection bias.
In other recent work, the FCI algorithm has been extended to incorporate prior knowledge regarding exogeneity and unconfoundedness of context variables (of the same type that we used for modeling the treatment variable in randomized controlled trials) \cite{Mooij++_JMLR_20}.
Such an extension with exogenous input nodes allows to generalize the idea of causal discovery in a randomized controlled trial to multiple treatment and outcome variables.
We also incorporate that extension here, and provide a `unified' extended FCI algorithm that allows for cycles (under certain assumptions), selection bias and exogenous input variables.

%We provide here a treatment of FCI that is unique in the sense that it is self-contained, it allows for cycles by assuming that the data was generated by a simple SCM, and it allows for the presence of exogenous input nodes.
%While the extended FCI of \cite{SMR1999} can deal with selection bias, in our treatment here we will assume that no selection bias is present.\footnote{Deriving a version of FCI that allows to deal with cycles, exogenous input nodes and selection bias is an open research problem.}

\subsection{Modeling selection bias}

We model selection bias as follows. 
\begin{Not}
We will assume the existence of a simple SCM $M = \lt \Sin, \Send^+, \Srnd, \CX, P, f \rt$, with endogenous variables $V^+ = V \dcup S$.\footnote{Alternatively, one could assume endogenous variables $V^+ = V \dcup S \dcup L$ where $S$ are used as selection variables. We could then start by marginalizing out the variables in $L$ to arrive at the setting that is our starting point here.}
We assume that only the endogenous variables in $V$ are observed, as well as the exogenous input variables $J$.
The latent variables in $S$ have the role of \emph{latent selection variables}.
In other words, we will assume that the data is distributed according to the marginal Markov kernel of $M$ on $V$ after conditioning on $S$:
$$P_M(X_V \mid X_S \in \xi_S, \doit(X_J))$$
for some measurable set $\xi_S \subseteq \CX_S$.\footnote{Note that we did not include the exogenous random variables here; we are treating them as if they were unobserved. If some (or all) exogenous random variables are observed, then we have three options: (i) include them, but ignore the additional background knowledge that they are independent (that is, treating them as if they were endogenous), (ii) make endogenous copies and include those; (iii) include them and make use of the additional background knowledge that they are independent. Here, we chose for the second option.}
\end{Not}
This models a process that \emph{filters} the data according to the value of $X_S$ (like in a rejection sampler).\footnote{It is helpful to think about such a filtering step as an intervention on the population level, but may lead to confusion when interpreted as an intervention on the individual level.}
One practical example of such a filtering process leading to selection bias is the following.
\begin{Eg}
  After the exam, the teacher hands out evaluation forms to the students and asks them to evaluate the course.
  When analyzing this evaluation, the teacher needs to be aware of possible selection bias.
  For example, if the students that were most dissatisfied with the course already dropped out earlier on and never filled in the evaluation form, the observed evaluations may not be representative of the opinions of all the students that started the course.
  Here, the selection variable may be the binary indicator `student filled in the evaluation form'.
\end{Eg}
  
Our goal in the rest of this chapter will be to deduce as much as possible about the graph $G_{V^+|J}(M)$ from the Markov kernel $P_{M}(X_V \mid X_S \in \xi_S, \doit(X_J))$. 
FCI is an example of a \emph{constraint-based} causal discovery algorithm, which means that it only utilizes the conditional independence information of this Markov kernel.
\Joris{I believe I need to extend the following to a ``data generating process'', a combination of an SCM and a set of selection events.}

\begin{Prp}
  Given a Markov kernel $K(X_V \mid X_J)$.
  Suppose there exists a simple SCM $M = \lt \Sin, \Send^+, \Srnd, \CX, P, f \rt$, with endogenous variables $V^+ = V \dcup S$,
  such that
  $$K(X_V \mid X_J) = P_M(X_V \mid X_S \in \xi_S, \doit(X_J))$$
  for some measurable set $\xi_S \subseteq \CX_S$.
  Then there exists a simple SCM $\tilde{M} = \lt \Sin, \tilde\Send^+, \Srnd, \CX, P, \tilde{f} \rt$, with endogenous variables $\tilde{V}^+ = V \dcup \{\tilde{s}\}$,
  such that
  $$P_M(X_V \mid X_S \in \xi_S, \doit(X_J)) = P_{\tilde M}(X_V \mid X_{\tilde s} = 1, \doit(X_J))$$
  and $\Anc_{G(M)}(S) = \Anc_{G(\tilde{M})}(\tilde{s})$.
\end{Prp}
\begin{proof}
  Define $\hat{M} = \lt \Sin, \hat\Send^+, \Srnd, \CX, P, \hat{f} \rt$, with endogenous variables $\hat{V}^+ = V \dcup S \dcup \{\tilde{s}\}$,
  and $\hat{f}_V := f_V$, $\hat{f}_S := f_S$, $\hat{f}_{\tilde{s}}(x) := \ind_{\xi_S}(x_{S})$.
  Now take $\tilde{M} = \hat{M}^{\sm S}$.
\end{proof}

\begin{Not}
We will assume the existence of a simple SCM $M = \lt \Sin, \Send^+, \Srnd, \CX, P, f \rt$, with endogenous variables $V^+ = V \dcup S$.\footnote{Alternatively, one could assume endogenous variables $V^+ = V \dcup S \dcup L$ where $S$ are used as selection variables. We could then start by marginalizing out the variables in $L$ to arrive at the setting that is our starting point here.}
We assume that only the endogenous variables in $V$ are observed, as well as the exogenous input variables $J$.
The latent variables in $S$ have the role of \emph{latent selection variables}.
In other words, we will assume that the data is distributed according to the marginal Markov kernel of $M$ on $V$ after conditioning on $S$:
$$P_M(X_V \mid X_S \in \xi_S, \doit(X_J))$$
for some measurable set $\xi_S \subseteq \CX_S$.\footnote{Note that we did not include the exogenous random variables here; we are treating them as if they were unobserved. If some (or all) exogenous random variables are observed, then we have three options: (i) include them, but ignore the additional background knowledge that they are independent (that is, treating them as if they were endogenous), (ii) make endogenous copies and include those; (iii) include them and make use of the additional background knowledge that they are independent. Here, we chose for the second option.}
\end{Not}

\subsection{Inducing walks}

The key notion in this chapter is that of \emph{$\sigma$-inducing walks}.
We define $\sigma$-inducing walk as a generalization to CDMGs of the notion of inducing path \cite{VermaPearl1990} in DAGs. %(primitive inducing path in \cite{RichardsonSpirtes02}).
\begin{Def}\label{def:sigma-inducing_path}
Let $G$ be a CDMG with input nodes $J$, output nodes $V^+ = V \dcup S$ and let $i,j \in V \dcup J$ be distinct nodes.
A walk $\pi$ in $G$ between $i$ and $j$ is called \emph{$\sigma$-inducing given $S$} if  each collider on $\pi$ is in $\Anc_{G}(\{i,j\}\cup S)$, and each non-endpoint non-collider on $\pi$ is unblockable. 
If it is a path, it is called a \emph{$\sigma$-inducing path given $S$} between $i$ and $j$.\footnote{In most of the literature, the graph is assumed to be acyclic, and then the notion is referred to simply as ``inducing''.}
\end{Def}
If two nodes are adjacent in $G$, any edge connecting the two is a $\sigma$-inducing walk (path) given $S$ between them, for any $S$.
Figure~\ref{fig:sigma-inducing_path} shows some simple nontrivial examples of $\sigma$-inducing paths.
\begin{Lem}\label{lem:sigma-inducing_path_in_scc}
Let $G$ be a CDMG with input nodes $J$, output nodes $V^+ = V \dcup S$ and let $i,j \in V \dcup J$ be distinct nodes.
If $i \in \Sc_G(j)$ then there exists a $\sigma$-inducing path given $S$ in $G$ between $i$ and $j$.
\end{Lem}
\begin{proof}
There exists a directed path in $G$ from $i$ to $j$ that is entirely contained in $\Sc_G(j)$, and therefore all its non-endpoint nodes are unblockable non-colliders.
\end{proof}

\begin{figure}\centering
  \begin{tikzpicture}
    \begin{scope}
      \node at (-4,0) {(a)};
      \node[ndout] (X1) at (-3,0) {$i$};
      \node[ndout] (X2) at (-1.5,0) {$j$};
      \node[ndout] (X3) at (0,0) {$k$};
      \draw[tuh] (X1) edge (X2);
      \draw[tuh] (X2) edge (X3);
      \draw[huh,bend left] (X2) edge (X3);
    \end{scope}
    \begin{scope}[xshift=5.5cm]
      \node at (-4,0) {(b)};
      \node[ndout] (X1) at (-3,0) {$i$};
      \node[ndout] (X2) at (-1.5,0) {$j$};
      \node[ndout] (X3) at (0,0) {$k$};
      \draw[tuh] (X1) edge (X2);
      \draw[tuh] (X2) edge (X3);
      \draw[tuh,bend left] (X3) edge (X2);
    \end{scope}
    \begin{scope}[xshift=11cm]
      \node at (-4,0) {(c)};
      \node[ndout] (X1) at (-3,0) {$i$};
      \node[ndout] (X2) at (-1.5,0) {$j$};
      \node[ndout] (X3) at (0,0) {$k$};
      \node[ndout] (X4) at (-1.5,-1) {$l$};
      \draw[tuh] (X1) edge (X2);
      \draw[tuh] (X2) edge (X3);
      \draw[tuh] (X3) edge (X4);
      \draw[tuh] (X4) edge (X1);
    \end{scope}
    \begin{scope}[yshift=-2cm]
      \node at (-4,0) {(d)};
      \node[ndint] (X1) at (-3,0) {$j$};
      \node[ndint] (X2) at (-1.5,0) {$i$};
      \node[ndout] (X3) at (0,0) {$k$};
      \draw[tuh] (X2) edge (X3);
    \end{scope}
    \begin{scope}[yshift=-2cm,xshift=5.5cm]
      \node at (-4,0) {(e)};
      \node[ndout] (X1) at (-3,0) {$i$};
      \node[ndout] (X2) at (-1,0) {$k$};
      \node[ndout] (X3) at (-2,-1) {$s$};
      \draw[tuh] (X1) edge (X3);
      \draw[tuh] (X2) edge (X3);
    \end{scope}
    \begin{scope}[yshift=-2cm,xshift=11cm]
      \node at (-4,0) {(f)};
      \node[ndint] (X1) at (-3,0) {$i$};
      \node[ndout] (X2) at (-1.5,0) {$l$};
      \node[ndout] (X3) at (-2.25,-1) {$s$};
      \node[ndout] (X4) at (0,0) {$k$};
      \draw[tuh] (X1) edge (X3);
      \draw[tuh,bend left] (X2) edge (X3);
      \draw[tuh,bend left] (X3) edge (X2);
      \draw[tuh] (X4) edge (X2);
    \end{scope}
  \end{tikzpicture}
  \caption{Examples of non-trivial $\sigma$-inducing paths in CDMGs. (a) The path $i \tuh j \huh k$ is a $\sigma$-inducing path between $i$ and $k$. (b) The path $i \tuh j \tuh k$ is a $\sigma$-inducing path between $i$ and $k$. (c) The path $i \tuh j \tuh k$ is a $\sigma$-inducing path between $i$ and $k$. The nodes $i$ and $k$ cannot be $\sigma$-separated by any subset not containing $i,k$ (indeed, $i \nPerp^\sigma k$ and $i \nPerp^\sigma k \mid j$ in all three graphs, and additionally $i \nPerp^\sigma k \mid l$ and $i \nPerp^\sigma k \mid \{j,l\}$ in the graph in (c)). (d) The path $i \tuh k$ is $\sigma$-inducing, while there is no $\sigma$-inducing path between $i$ and $j$, nor between $j$ and $k$. (e) The path $i \tuh s \hut k$ is $\sigma$-inducing given $\{s\}$. (f) The path $i \tuh s \tuh l \hut k$ is $\sigma$-inducing given $\{s\}$.\label{fig:sigma-inducing_path}}
\end{figure}

The notion of $\sigma$-inducing walk has the following important properties.
\begin{Prp}\label{prp:sigma-inducing_path_properties}
  Let $G$ be a CDMG with input nodes $J$, output nodes $V^+ = V \dcup S$ and let $i \in V$ (but $i \notin J$) and $j \in V \dcup J$ be distinct nodes. 
  Then the following are equivalent:
  \begin{enumerate}[label=(\roman*)]
    \item There is a $\sigma$-inducing path given $S$ in $G$ between $i$ and $j$;\label{prp:sigma-inducing_path_properties2:i}
    \item There is a $\sigma$-inducing walk given $S$ in $G$ between $i$ and $j$;\label{prp:sigma-inducing_path_properties2:ii}
    \item $i \nPerp^\sigma_G j \given S \cup Z$ for all $Z \subseteq (V \cup J) \sm \{i,j\}$;\label{prp:sigma-inducing_path_properties2:iii}
    \item $i \nPerp^\sigma_G j \given S \cup Z$ for $Z = (\Anc_G(\{i,j\} \cup S) \cup J) \sm \{i,j\}$.\label{prp:sigma-inducing_path_properties2:iv}
  \end{enumerate}
\end{Prp}
\begin{proof}
  The proof is similar to that of Theorem 4.2 in \cite{RichardsonSpirtes02}.

  \ref{prp:sigma-inducing_path_properties2:i} $\implies$ \ref{prp:sigma-inducing_path_properties2:ii} is trivial.

  \ref{prp:sigma-inducing_path_properties2:ii} $\implies$ \ref{prp:sigma-inducing_path_properties2:iii}: Assume the existence of a $\sigma$-inducing walk given $S$ between $i$ and $j$ in $G$. 
  Let $Z \subseteq (V \cup J) \sm\{i,j\}$. 
  Consider all walks in $G$ between $i$ and $j$ with the property that all colliders on it are in $\Anc_{G}({\{i,j\} \cup S \cup Z)}$, and each non-endpoint non-collider on it is not in $S \cup Z$ or is unblockable.
  Such walks exist, since the $\sigma$-inducing walk is one. 
  Let $\mu$ be such a walk with a minimal number of colliders.
  We show that all colliders on $\mu$ must be in $\Anc_{G}(S \cup Z)$. 
  Suppose on the contrary the existence of a collider $k$ on $\mu$ that is not ancestor of $S \cup Z$.
  It is either ancestor of $i$ or of $j$, by assumption.
  If $j \in J$, it cannot be ancestor of $j$, and hence must be ancestor of $i$.
  Otherwise, we can assume it to be ancestor of $i$ without loss of generality.
  Then there is a directed path $\pi$ from $k$ to $i$ in $G$ that does not pass through any node of $S \cup Z$.
%  Let $k$ be the vertex closest to $j$ on $\mu$ that is also on $\pi$. 
  Then the subwalk of $\mu$ between $k$ and $j$ can be concatenated with the directed path $\pi$
  into a walk between $i$ and $j$ that has the property, but has fewer colliders than $\mu$: a contradiction. 
  Therefore, $\mu$ is $\sigma$-open given $S \cup Z$.
  %can be extended to a walk that is $\sigma$-open given $Z$, by replacing each collider $k$ on
  %it by a walk $k \to \dots \to z \ot \dots \ot k$ for some closest descendant $z \in Z$ of $k$ (with possibly $z=k$).
  Hence, $i$ and $j$ are $\sigma$-connected given $S \cup Z$.

  \ref{prp:sigma-inducing_path_properties2:iii} $\implies$ \ref{prp:sigma-inducing_path_properties2:iv} is trivial.

  \ref{prp:sigma-inducing_path_properties2:iv} $\implies$ \ref{prp:sigma-inducing_path_properties2:i}:
  Suppose that %$i$ and $j$ are $\sigma'$-connected given any $Z \subseteq V \sm \{i,j\}$.
  %In particular, 
  $i$ and $j$ are $\sigma$-connected given $Z = (\Anc_{G}(\{i,j\} \cup S) \cup J) \sm \{i,j\}$. 
  Let $\pi$ be a path between $i$ and $\{j\} \cup J$ that is $\sigma$-open given $Z$. % and contains $i$ and $j$ only once.
  The end nodes of $\pi$ must be $i$ and $j$, because $J \sm \{i,j\} \subseteq Z$.
  We show that $\pi$ must be a $\sigma$-inducing path given $S$. 
  First, all colliders on $\pi$ are in $\Anc_G(Z)$, but not in $J$, and hence in $\Anc_{G}(\{i,j\} \cup S)$.
  Second, let $k$ be any non-endpoint non-collider on $\pi$.
  Then there must be a directed subpath of $\pi$ starting at $k$ that ends either at the first collider on $\pi$ next to $k$ or at an end node of $\pi$, and hence $k$ must be in $Z$. 
  Since $\pi$ is $\sigma$-open given $Z$, $k$ must be unblockable.
  Hence, all non-endpoint non-colliders on $\pi$ must be unblockable.
\end{proof}
In words: there is a $\sigma$-inducing path between two nodes in a CDMG (provided they are not both input nodes) if and only if the two nodes cannot be $\sigma$-separated by any subset of the other nodes.
%A similar property holds for $d$-inducing walks and paths, but with $\sigma$-separation replaced by $d$-separation in points (iii) and (iv).
\begin{Rem}
Two input nodes cannot be $\sigma$-separated by some subset of other nodes.
Indeed, the trivial path $j$ is $\sigma$-open as long as we don't condition on $j$ itself.
On the other hand, there may, or may not, be a $\sigma$-inducing path given $S$ between two input nodes. 
For example, in the CDMG $j_1 \tuh i \hut j_2$ with $j_1,j_2 \in J$ the path $j_1 \tuh i \hut j_2$
is $\sigma$-inducing given $\{i\}$ but not $\sigma$-inducing given $\emptyset$.
This is why Proposition~\ref{prp:sigma-inducing_path_properties} does not consider the case $i,j \in J$.
\end{Rem}

The orientations of the outermost edges on a $\sigma$-inducing path contain important information about ancestral relations.
\begin{Lem}\label{lem:sigma-inducing_path_out-of}
  Let $G$ be a CDMG with input nodes $J$, output nodes $V^+ = V \dcup S$ and let $i,j \in V \dcup J$ be distinct.
  If there exists a $\sigma$-inducing path given $S$ between $i$ and $j$ in $G$,
  %and all $\sigma$-inducing paths in $G$ between $i$ and $j$ are out of $i$, then $i \in \Anc_G(j)$.
  and all $\sigma$-inducing paths given $S$ in $G$ between $i$ and $j$ are out of $j$, then $j \in \Anc_G(\{i\} \cup S)$.
\end{Lem}
\begin{proof}
  Let $\mu$ be a $\sigma$-inducing path given $S$ between $i$ and $j$ in $G$. 
  It must be of the form $i \cdots l \hut j$ (with possibly $l = i$). 
  First we show that $l$ cannot be in $\Sc_G(j)$.
  If $l \in \Sc_G(j)$, then let $\pi$ be a directed path in $G$ from $l$ to $j$ that is entirely contained in $\Sc_G(j)$.
  Let $m$ be the node on $\mu$ closest to $i$ that is also on $\pi$ (possibly $m = l$).
  The subpath of $\pi$ between $j$ and $m$ can be concatenated with the subpath of $\mu$ between $m$ and $i$ into a walk between $j$ and $i$.
  This must be a $\sigma$-inducing path given $S$ between $i$ and $j$ that is into $j$ by construction: contradiction.
  Hence $l$ cannot be in $\Sc_G(j)$.

  If $\mu$ is a directed path all the way to $i$, then clearly, $j \in \Anc_G(\{i\} \cup S)$.
  Otherwise, it must contain a collider.
  Let $k$ be the collider on $\mu$ closest to $j$.
  $k$ must be ancestor of $i$ or $j$ or $S$.
  In the first and third cases, clearly $j \in \Anc_G(\{i\} \cup S)$.
  In the second case, all nodes on the subpath of $\mu$ between $j$ and $k$ must be in $\Sc_G(j)$, a contradiction.
\end{proof}

\begin{Lem}\label{lem:sigma-inducing_path_into_into}
  Let $G$ be a CDMG with input nodes $J$, output nodes $V^+ = V \dcup S$ and let $i,j \in V \dcup J$ be distinct.
  If there exists a $\sigma$-inducing path given $S$ between $i$ and $j$ in $G$ into $j$, and $i \notin \Anc_G(\{j\} \cup S)$, then there exists a $\sigma$-inducing path given $S$ between $i$ and $j$ in $G$ that is both into $i$ and into $j$.
\end{Lem}
\begin{proof}
  Let $\mu$ be a $\sigma$-inducing path given $S$ between $i$ and $j$ in $G$ into $j$.
  This rules out $j \in J$.
  If $\mu$ is into $i$, we are done.
  Therefore, suppose it is of the form $i \tuh \cdots \suh j$. 
  It cannot be a directed path, since $i \notin \Anc_G(\{j\} \cup S)$.
  Therefore, there must be a collider $k$ on $\mu$ such that $\mu$ is of the form $i \tuh \dots \tuh k \hus \cdots \suh j$
  (with the subpath between $i$ and $k$ directed).
  Then $k \in \Anc_G(\{i\})$ (sic!), and hence all nodes on $\mu$ between $i$ and $k$ must be in $\Sc_G(i)$.
  Let $\pi$ be a directed path in $G$ from $k$ to $i$ that is entirely contained in $\Sc_G(i)$.
  Let $l$ be the node on $\mu$ closest to $j$ that is also on $\pi$ (possibly $l = k$).
  Then $l \ne j$, because otherwise $j \in \Sc_G(i)$, contradicting $i \notin \Anc_G(\{j\} \cup S)$.
  The non-trivial subpath of $\pi$ between $i$ and $l$ can be concatenated with the non-trivial subpath of $\mu$ between $l$ and $j$ into a walk between $i$ and $j$.
  This must be a $\sigma$-inducing path given $S$ between $i$ and $j$ that is both into $i$ and into $j$.
\end{proof}

\begin{Lem}\label{lem:inducing_path_into_scc}
Let $G$ be a CDMG with input nodes $J$, output nodes $V^+ = V \dcup S$ and let $i,j \in V \dcup J$ be distinct.
If there is a $\sigma$-inducing path in $G$ given $S$ between $i$ and $j$ that is into $j$, and $k \in \Sc_G(j)$, then there is a $\sigma$-inducing path in $G$ given $S$ between $i$ and $k$ that is into $k$.
\end{Lem}
\begin{proof}
The $\sigma$-inducing path given $S$ in $G$ between $i$ into $j$ can be extended by a directed path within $\Sc_G(k)$ to a $\sigma$-inducing walk given $S$ in $G$ between $i$ and $k$ that is into $k$.
\end{proof}

\begin{Lem}\label{lem:inducing_path_out_of_input}
  Let $G$ be a CDMG with input nodes $J$, output nodes $V^+ = V \dcup S$ and let $i,j \in V \dcup J$ be distinct.
  If there is a $\sigma$-inducing path in $G$ given $S$ between $j \in J$ and $v \in V$, then $j \in \Anc_G(\{v\} \cup S)$. 
\end{Lem}
\begin{proof}
  The $\sigma$-inducing path in $G$ must start out of $j$ by definition. 
  If it is a directed path from $j$ all the way to $v$, then clearly $j \in \Anc_G(v) \subseteq \Anc_G(\{v\} \cup S)$.
  Otherwise, it will contain a directed path from $j$ to a collider.
  That collider must be in $\Anc_G(\{j,v\} \cup S)$, but since it cannot be in $\Anc_G(j)$, it has to be in $\Anc_G(\{v\} \cup S)$.
  This means that $j \in \Anc_G(\{v\} \cup S)$.
\end{proof}

\subsection{Partial Ancestral Graphs}

It is often convenient when performing causal reasoning to be able to represent a set of CDMGs in a compact way. 
For this purpose, \emph{partial ancestral graphs (PAGs)} have been introduced \cite{SMR1999,Zhang2006}.
In order to deal with possible cycles (in simple SCMs), selection bias and exogenous input nodes, we extend the definition to $\sigma$-PAGs.
%Here we will assume no selection bias for simplicity, which allows us to restrict ourselves to discussing \emph{directed} PAGs (henceforth abbreviated as DPAGs).
%We will extend these objects to allow for input nodes, in which case we will refer to them as \emph{conditional directed} PAGs ($\sigma$-PAGs).

$\sigma$-PAGs have two node types, input nodes and output nodes (just like CDMGs).
They also have multiple edge types. In addition to the three edge types for CDMGs
($\tuh$, $\hut$, $\huh$), there is an undirected edge ($\tut$) and there are five
edge types involving a circle edge mark: $\huo$, $\ouo$, $\ouh$, $\tuo$, $\out$. 
Each edge $i \sus j$ has two \emph{edge marks}, one at each node, with each edge mark
either a \emph{tail}, \emph{arrowhead} or \emph{circle}. For example, the directed edge $i \tuh j$ has a 
tail at $i$ and an arrowhead at $j$, while the bi-circle edge $i \ouo j$ has two circle edge marks.
All 9 possible combinations of edge marks can occur on an edge in a $\sigma$-PAG. We will make use
of the ``$\asterisk$'' symbol to denote any of the three edge marks. So the notation $i \sus j$
can stand for all 9 possible edge types between $i$ and $j$, whereas $i \hus j$ is shorthand for three
possible edge types, as are $i \tus j$ and $i \ous j$.
Edges of the form $i \hus j$ and $j \suh i$ are called \emph{into $i$}. 
Edges of the form $i \tus j$ and $j \sut i$ are called \emph{out of $i$}.

In order to define $\sigma$-PAGs, we extend the definitions of (directed) walks, (directed) paths and colliders to cover these new edge types.
\begin{Def}\label{def:more_walks}
  Let $H=\lt J,V,E \rt$ be a mixed graph with input nodes $J$, output nodes $V$ and edges $E$ of the types $\{\tuh,\hut,\huo,\huh,\ouo,\ouh,\tuo,\out,\tut\}$.\footnote{Formally, we no longer introduce separate sets to represent the edges of each type, but merge them into the single set $E$.}
  Let $v,w \in V \cup J$.
  \begin{enumerate}
    \item If there is an edge $v \sus w$ between $v$ and $w$ (of any type), we call $v$ and $w$ \emph{adjacent} in $H$.
    \item Define $\Adj_H(v)$ to be the nodes adjacent to $v$ in $H$.
    \item A triple of distinct nodes $\lt a,b,c \rt$ in $H$ form a \emph{triangle} if each pair of nodes in the triple is adjacent in $H$.
    \item A triple of distinct nodes $\lt a,b,c \rt$ in $H$ is called \emph{unshielded} if $b$ is adjacent to both $a$ and $c$ in $H$, but $a$ is not adjacent to $c$ in $H$.
    \item A \emph{walk between $v$ and $w$} in $H$ is a finite alternating sequence of nodes and edges 
        \[v=v_0, a_0, v_1, \dots v_{n-1}, a_{n-1}, v_n=w\] 
        in $H$ for some $n \ge 0$, i.e.\ such that for every $k=0,\dots,n-1$ 
        we have that $v_{k},v_{k+1} \in V \cup J$ and $a_k = v_k \sus v_{k+1} \in E$, and with end nodes $v_0=v$ and $v_n=w$.
    \item A walk is called a \emph{path} if no node occurs more than once on the walk.
    \item A \emph{directed walk} (path) from $v$ to $w$ in $H$ is a walk (path) of the form:
        \[v=v_0 \tuh v_1 \tuh  \cdots \tuh v_{n-1} \tuh v_n=w,\]
        for some $n \ge 0$.
    \item A path $v_0 \sus \dots \sus v_n$ in $H$ is called a \emph{possibly directed path from $v_0$ to $v_n$} if for each $i=1,\dots,n$, the edge $v_{i-1} \sus v_i$ is not into $v_{i-1}$ and is not out of $v_i$ (i.e., each edge must be of the form $v_{i-1} \ouo v_i$, $v_{i-1} \ouh v_i$, $v_{i-1} \tuo v_i$ or $v_{i-1} \tuh v_i$).
    \item A \emph{directed cycle} is a directed walk $v \tuh \dots \tuh w$, concatenated with
      the directed edge $w \tuh v$.
    \item An \emph{almost directed cycle} is a directed walk $v \tuh \dots \tuh w$, concatenated
      with the bidirected edge $w \huh v$.
    \item A path of the form $v_0 \ouo \dots \ouo v_n$ (with each edge of the form $v_i \ouo v_{i+1}$) is called a \emph{circle path}.
    \item A path $v_0 \sus \dots \sus v_n$ in $H$ is called \emph{uncovered} if every subsequent triple $\lt v_{k-1}, v_k, v_{k+1} \rt$ for $1 < k < n$ is unshielded.
    \item A triple of consecutive nodes $v_{k-1} \sus v_k \sus v_{k+1}$ on a walk in $H$ is called \emph{definite collider}
      if it is of the form $v_{k-1} \suh v_k \hus v_{k+1}$.
    \item A triple of consecutive nodes $v_{k-1} \sus v_k \sus v_{k+1}$ on a walk in $H$ is called a \emph{definite non-collider}
      if it is of the form $v_{k-1} \sus v_k \tus v_{k+1}$ or $v_{k-1} \sut v_k \sus v_{k+1}$.
      Furthermore, we also refer to the end nodes $v_0$ and $v_n$ of a walk between $v_0$ and $v_n$ in $H$ as \emph{definite non-colliders}.
    \item If there is a directed walk from $v$ to $w$ in $H$ then we say that \emph{$v$ is ancestor of $w$ in $H$}, 
      and we write $v \in \Anc_H(w)$. For $W \subseteq V \cup J$, we define $\Anc_H(W) := \bigcup_{w\in W} \Anc_H(w)$.
    \item If there is a directed walk from $v$ to $w$ in $H$ then we say that \emph{$w$ is descendant of $v$ in $H$}, 
      and we write $w \in \Desc_H(v)$. For $W \subseteq V \cup J$, we define $\Desc_H(W) := \bigcup_{w\in W} \Desc_H(w)$.
    \item A walk between $v,w \in V \cup J$ (with $v \ne w$) in $H$ is called \emph{definitely inducing} if every non-endpoint node is a definite collider in $\Anc_H(\{v,w\})$.
    \item A walk between $v$ and $w$ in $H$ is called \emph{definitely open given $Z \subseteq V \cup J$} if
      \begin{enumerate}
        \item every node on the walk is a definite collider or definite non-collider, and
        \item every definite collider is in $\Anc_H(Z)$, and
        \item every definite non-collider is not in $Z$.
      \end{enumerate}
  \end{enumerate}
\end{Def}

We can now define:\footnote{We have incorporated two extensions of the usual definition of PAG \cite{Zhang2006}: we allow for input 
nodes, and we have weakened the condition of being ancestral to $\sigma$-ancestral. A mixed graph is called \emph{ancestral} if it
has no directed cycles, no almost directed cycles, and no triples of the form $i \suh j \tut k$.}
\begin{Def}\label{def:sigmapag}
%  We call a graph $G = \lt V,E\rt$ with nodes $V$ and edges $E$ between ordered node pairs of the types $\{\ttt,\ttc,\to,\ot,\otc,\oto,\ctt,\ctc,\cto\}$ a \emph{partial ancestral graph (PAG)} if
  %It is called a \emph{directed partial ancestral graph (DPAG)} if it only has edges of the types $\{\to,\ot,\oto,\otc,\cto,\ctc\}$.
  A mixed graph $H = \lt J,V,E\rt$ with input nodes $J$, output nodes $V$ and edges $E$ of the types $\{\tuh,\hut,\huo,\huh,\ouo,\ouh,\tuo,\out,\tut\}$ is called a \emph{partial $\sigma$-ancestral graph ($\sigma$-PAG)} if all of the following conditions hold:
  \begin{enumerate}
    \item Between any two distinct nodes there is at most one edge, and there are no edges between a node and itself;
    \item No input node is adjacent to any other input node;
    \item The graph contains no directed cycles, no almost directed cycles, and no unshielded triple of the form $i \suh j \tut k$ (``$\sigma$-ancestral'');
    \item There is no definitely inducing path between any two distinct non-adjacent nodes (``maximal'');
%    \item If there is an edge $j \sus v$ with $j \in J, v \in V$ then it must be of the form $j \tus v$.
  \end{enumerate}
\end{Def}

$\sigma$-PAGs are used to represent a set of CDMGs as follows.
\begin{Def}\label{def:sigmapag_represents_cdmg}
  Let $H = \lt J,V,E\rt$ be a mixed graph with input nodes $J$, output nodes $V$ and edges $E$ of the types $\{\tuh,\hut,\huo,\huh,\ouo,\ouh,\tuo,\out,\tut\}$.
  Let $G$ be a CDMG with input nodes $J$ and output nodes $V^+ = V \dcup S$.
  We say that \emph{$H$ represents $G$ given $S$} if all of the following hold:
  \begin{enumerate}
    \item Between any two distinct nodes in $H$ there is at most one edge, and there are no edges between a node and itself;
    \item Two distinct nodes $i,j \in V \cup J$ are adjacent in $H$ if and only if $\{i,j\} \not\subseteq J$ and there is a $\sigma$-inducing path between $i$ and $j$ given $S$ in $G$;
    \item If $i \hus j$ in $H$ then $i \notin \Anc_{G}(\{j\} \cup S)$;
    \item If $i \tus j$ in $H$ then $i \in \Anc_{G}(\{j\} \cup S)$;
%    \item If $j \sus v$ in $H$ with $j \in J, v \in V$ then it must be of the form $j \tus v$.
  \end{enumerate}
\end{Def}
Hence, adjacencies represent $\sigma$-inducing paths, arrowheads represent non-ancestorship (of the adjacent node or $S$), and tails represent ancestorship (of the adjacent node or $S$).
Some examples are given in Figure~\ref{fig:sigmapags}. 
Note in particular that a directed edge in a $\sigma$-PAG does not necessarily imply a direct causal relationship (for example, the edge $i \tuh k$ in Figure~\ref{fig:sigmapags}(a)).
Figure~\ref{fig:nonmaximal_ancestral_graph} provides an example of a mixed graph that satisfies all conditions of a $\sigma$-PAG except the maximality.

\begin{figure}\centering
  \begin{tikzpicture}
    \begin{scope}
      \node at (-4,0) {(a)};
      \node[ndout] (X1) at (-3,0) {$i$};
      \node[ndout] (X2) at (-1.5,0) {$j$};
      \node[ndout] (X3) at (0,0) {$k$};
      \draw[tuh] (X1) edge (X2);
      \draw[tuh] (X2) edge (X3);
%      \draw[biarr,bend left] (X2) edge (X3);
      \draw[tuh,bend right] (X1) edge (X3);
    \end{scope}
    \begin{scope}[xshift=5.5cm]
      \node at (-4,0) {(b)};
      \node[ndout] (X1) at (-3,0) {$i$};
      \node[ndout] (X2) at (-1.5,0) {$j$};
      \node[ndout] (X3) at (0,0) {$k$};
      \draw[tuh] (X1) edge (X2);
      \draw[ouo] (X2) edge (X3);
%      \draw[biarr,bend left] (X2) edge (X3);
      \draw[tuh,bend right] (X1) edge (X3);
    \end{scope}
    \begin{scope}[xshift=11cm]
      \node at (-4,0) {(c)};
      \node[ndout] (X1) at (-3,0) {$i$};
      \node[ndout] (X2) at (-1.5,0) {$j$};
      \node[ndout] (X3) at (0,0) {$k$};
      \draw[tut] (X1) edge (X2);
      \draw[tut] (X2) edge (X3);
%      \draw[biarr,bend left] (X2) edge (X3);
      \draw[tut,bend right] (X1) edge (X3);
    \end{scope}
    \begin{scope}[yshift=-2cm]
      \node at (-4,0) {(d)};
      \node[ndint] (X1) at (-3,0) {$j$};
      \node[ndint] (X2) at (-1.5,0) {$i$};
      \node[ndout] (X3) at (0,0) {$k$};
      \draw[tuh] (X2) edge (X3);
    \end{scope}
    \begin{scope}[yshift=-2cm,xshift=5.5cm]
      \node at (-4,0) {(e)};
      \node[ndout] (X1) at (-3,0) {$i$};
      \node[ndout] (X2) at (-1.5,0) {$k$};
      \draw[tut] (X1) edge (X2);
    \end{scope}
    \begin{scope}[yshift=-2cm,xshift=11cm]
      \node at (-4,0) {(f)};
      \node[ndout] (X1) at (-3,0) {$i$};
      \node[ndout] (X2) at (-1.5,0) {$l$};
      \node[ndout] (X3) at (0,0) {$k$};
      \draw[ouo] (X1) edge (X2);
      \draw[ouo] (X2) edge (X3);
      \draw[ouo,bend left] (X3) edge (X1);
    \end{scope}
    \begin{scope}[yshift=-4cm]
      \node at (-4,0) {(g)};
      \node[ndout] (X1) at (-3,0) {$X_1$};
      \node[ndout] (X2) at (-1,0) {$X_2$};
      \node[ndout] (X3) at (-2,-1) {$X_3$};
      \node[ndout] (X4) at (-2,-2.5) {$X_4$};
      \draw[ouh] (X1) edge (X3);
  %    \draw[biarr,bend right] (X1) edge (X3);
      \draw[ouh] (X2) edge (X3);
  %    \draw[biarr,bend left] (X2) edge (X3);
      \draw[tuh] (X3) edge (X4);
    \end{scope}
    \begin{scope}[yshift=-4cm,xshift=5.5cm]
      \node at (-4,0) {(h)};
      \node[ndout] (X1) at (-3,0) {$X_1$};
      \node[ndout] (X2) at (-1.5,0) {$X_2$};
      \node[ndout] (X3) at (0,0) {$X_3$};
      \draw[ouo] (X1) edge (X2);
      \draw[ouo] (X2) edge (X3);
    \end{scope}
    \begin{scope}[yshift=-4cm,xshift=11cm]
      \node at (-4,0) {(i)};
      \node[ndout] (X1) at (-3,0) {$X_1$};
      \node[ndout] (X2) at (-1.5,0) {$X_2$};
      \node[ndout] (X3) at (0,0) {$X_3$};
      \draw[ouh] (X1) edge (X2);
      \draw[tuh] (X2) edge (X3);
    \end{scope}
  \end{tikzpicture}
  \caption{Various example $\sigma$-PAGs, representing respectively the CDMGs:
  (a) in Figure~\ref{fig:sigma-inducing_path}(a);
  (b) in Figures~\ref{fig:sigma-inducing_path}(a--b);
  (c) in Figures~\ref{fig:sigma-inducing_path}(a--b);
  (d) in Figure~\ref{fig:sigma-inducing_path}(d);
  (e) in Figure~\ref{fig:sigma-inducing_path}(e);
  (f) in Figure~\ref{fig:sigma-inducing_path}(f);
  (g) of all Y-structures in Figure~\ref{fig:Ystruct};
  (h) of all LCD structures in Figure~\ref{fig:LCD};
  (i) of all LCD structures in Figure~\ref{fig:LCD}.\label{fig:sigmapags}}
\end{figure}

\begin{figure}\centering
  \begin{tikzpicture}
    \begin{scope}
%      \node at (-2,0) {(d)};
      \node[ndout] (X1) at (-1,0) {$i$};
      \node[ndout] (X2) at (-1,2) {$j$};
      \node[ndout] (X3) at (1,2) {$k$};
      \node[ndout] (X4) at (1,0) {$l$};
      \draw[huh] (X1) edge (X2);
      \draw[huh] (X2) edge (X3);
      \draw[huh] (X3) edge (X4);
      \draw[tuh] (X2) edge (X4);
      \draw[tuh] (X3) edge (X1);
    \end{scope}
  \end{tikzpicture}
  \caption{This mixed graph is not a valid $\sigma$-PAG, because it is not maximal: 
  it has a definitely inducing path $i \huh j \huh  k \huh l$ while $i$ and $l$ are non-adjacent.\label{fig:nonmaximal_ancestral_graph}}
\end{figure}

%The following property shows that we can no longer unambiguously read off ancestral relations from a $\sigma$-PAG that represents a CDMG in case of selection bias ($S \ne \emptyset$).
%The following lemma shows that the orientation of edges in a $\sigma$-PAG that represents a CDMG contains information on the orientation of corresponding $\sigma$-inducing paths in the CDMG.
The following elementary properties shows that certain important properties of the CDMG are encoded in a mixed graph that represents it.
\begin{Lem}\label{lem:mixed_graph_represents_CDMG}
  Let $H = \lt J,V,E\rt$ be a mixed graph with input nodes $J$, output nodes $V$ and edges $E$ of the types $\{\tuh,\hut,\huo,\huh,\ouo,\ouh,\tuo,\out,\tut\}$.
  Let $G$ be a CDMG with input nodes $J$ and output nodes $V^+ = V \dcup S$.
  Suppose that $H$ represents $G$ given $S$.
  Then for any two nodes $i,j$ in $H$: 
  \begin{enumerate}[label=(\roman*)]
    \item $i \in \Anc_{H}(j)$ implies $i \in \Anc_{G}(\{j\} \cup S)$.
    \item If $i \suh j$ in $H$, then there exists a $\sigma$-inducing path given $S$ in $G$ between $i$ and $j$ that is into $j$. \label{lem:sigma-inducing_path_into2_into}
    \item If $i \huh j$ in $H$, then there exists a $\sigma$-inducing path given $S$ in $G$ between $i$ and $j$ that is both into $i$ and into $j$. \label{lem:sigma-inducing_path_into2_intointo}
%    \item If $i \hus j$ in $H$, then there exists a $\sigma$-inducing path in $G$ between $i$ and $j$ that is into $i$.
  \end{enumerate}
\end{Lem}
\begin{proof}
  \begin{enumerate}[label=(\roman*)]
    \item Suppose $H$ contains a directed path $i = v_0 \tuh \dots \tuh v_n = j$.
      Note that for all $k = 0, \dots, n-1$, $v_k \notin \Anc_G(S)$ implies $v_k \in \Anc_G(v_{k+1})$. %and $v_{k+1} \notin \Anc_G(S)$.
      By induction, then, $v_0 \notin \Anc_G(S)$ implies $v_0 \in \Anc_G(v_n)$.
    \item There exists a $\sigma$-inducing path given $S$ between $i$ and $j$ in $G$ because $i$ and $j$ are adjacent in $H$ and $H$ represents $G$ given $S$.
      If all $\sigma$-inducing paths given $S$ between $i$ and $j$ in $G$ were out of $j$, then by Lemma~\ref{lem:sigma-inducing_path_out-of},
      $j \in \Anc_G(\{i\} \cup S)$, contradicting the orientation $i \suh j$ in $H$. 
      Therefore, there must be a $\sigma$-inducing path given $S$ between $i$ and $j$ in $G$ that is into $j$.
      This shows~\ref{lem:sigma-inducing_path_into2_into}.
    \item This follows from case \ref{lem:sigma-inducing_path_into2_into} in combination with Lemma~\ref{lem:sigma-inducing_path_into_into}.
  \end{enumerate}
\end{proof}

We will frequently use that every mixed graph (of a certain type) that represents a CDMG must be a valid $\sigma$-PAG.
\begin{Prp}\label{prp:sigmapag_represents_cdmg_is_valid}
  Let $H = \lt J,V,E\rt$ be a mixed graph with input nodes $J$, output nodes $V$ and edges $E$ of the types $\{\tuh,\hut,\huo,\huh,\ouo,\ouh,\tuo,\out,\tut\}$.
  Let $G$ be a CDMG with input nodes $J$ and output nodes $V^+ = V \dcup S$.
  If $H$ represents $G$ given $S$, then $H$ is a $\sigma$-PAG.
\end{Prp}
\begin{proof}
  By assumption: there is at most one edge between two distinct nodes; there are no edges between a node and itself;
  no input node is adjacent to any other input node. 
  %if there is an edge $j \sus v$ in $H$ with $j \in J, v \in V$ then it must be of the form $j \tus v$.
  
  We show that $H$ is $\sigma$-ancestral. 
  First, suppose that $H$ contained a directed path $i = v_0 \tuh \dots \tuh v_n = j$.
  By Lemma~\ref{lem:mixed_graph_represents_CDMG}, then $i \in \Anc_G(\{j\} \cup S)$.
  If $H$ contained an edge $j \suh i$, this would imply $i \notin \Anc_G(\{j\} \cup S)$, a contradiction. 
  Hence such edges cannot occur.
  This means that no directed cycles and no almost directed cycles occur in $H$.
  %The remaining statement to prove that $H$ is $\sigma$-ancestral follows by Lemma~\ref{lem:into_undirected_edge}.
%  Second, suppose that an unshielded triple of the form $i \suh j \tut k$ occured in $H$.
%  Since $j \in \Anc_G(\{k\} \cup S)$ but $j \notin \Anc_G(\{i\} \cup S)$, we must have $j \in \Anc_G(k)$.
%  Also, $k \in \Anc_G(\{j\} \cup S)$.
%  Since $j \in \Anc_G(k)$ and $j \notin \Anc_G(S)$, we must have $k \in \Anc_G(j)$.
%  But then $j \in \Sc_G(k)$, which is not possible according to Lemma~\ref{lem:induced_path_into_scc}.
%  Hence such unshielded triples cannot occur in $H$.
%  Together, this shows that $H$ is $\sigma$-ancestral.
  It remains to show that no unshielded triple of the form $i \suh j \tut k$ can occur in $H$.
  We prove this by contradiction.
  Since $j \in \Anc_G(\{k\} \cup S)$ but $j \notin \Anc_G(\{i\} \cup S)$, we must have $j \in \Anc_G(k)$.
  Also, $k \in \Anc_G(\{j\} \cup S)$.
  Since $j \in \Anc_G(k)$ and $j \notin \Anc_G(S)$, we must have $k \in \Anc_G(j)$.
  But then $j \in \Sc_G(k)$, which gives a contradiction with Lemma~\ref{lem:inducing_path_into_scc}.

  We continue to show that $H$ is maximal.
  Suppose there is a definitely inducing path $\mu$ in $H$ between two distinct nodes $u,v \in V \cup J$. 
  We first show that this implies that $\{u,v\} \not\subseteq J$.
  If $u,v \in J$, then a definitely inducing path between them cannot consist of a single edge, because all node pairs in $J$ are non-adjacent in $H$ by assumption.
  Hence it must be of the form $u \suh w_1 \dots w_n \hus v$ (with $n \ge 1$) with $w_1, w_n \in V$. 
  But since all nodes $w_1,\dots,w_n$ are definite colliders, none of them is in $\Anc_G(S)$.
  Since all nodes $w_1,\dots,w_n$ are in $\Anc_H(\{u,v\})$, by Lemma~\ref{lem:mixed_graph_represents_CDMG}, they must all be in $\Anc_G(\{u,v\})$.
  Since $G$ is a CDMG, they must all be in $\{u,v\} \subseteq J$, a contradiction.
  %Both $w_1$ and $w_n$ must be in $\Anc_H(\{u,v\})$, and hence $w_1, w_n \in \Anc_G(\{u,v\} \cup S)$ by Lemma~\ref{lem:mixed_graph_represents_CDMG}.
  %Both $w_1$ and $w_n$ are definite colliders, hence $w_1 \notin \Anc_G(\{u\} \cup S)$ and $w_n \notin \Anc_G(\{v\} \cup S)$.
  %Therefore, $w_1 \in \Anc_G(v)$ and $w_n \in \Anc_G(w)$.
  %This gives a contradiction with Lemma~\ref{lem:inducing_path_out_of_input}.
  Hence, there cannot be a definitely inducing path in $H$ between two nodes in $J$.

  Every edge $i \sus j$ on $\mu$ corresponds with a $\sigma$-inducing path $\pi_{ij}$ given $S$ in $G$ between $i$ and $j$.
  By Lemma~\ref{lem:mixed_graph_represents_CDMG}, these $\sigma$-inducing paths can be chosen to be into $i$ if the edge is $i \hus j$, into $j$ if the edge is $i \suh j$, and both into $i$ and $j$ if the edge is $i \huh j$. 
  Concatenate all $\pi_{ij}$ (following the edge ordering of $\mu$) into a walk $\pi$ in $G$ between $u$ and $v$.
  Every non-endpoint node on $\mu$ is a definite collider on $\mu$, by assumption.
  By construction, these nodes then become colliders on $\pi$.
  Since definite colliders on $\mu$ are in $\Anc_H(\{u,v\})$ by assumption, they are in $\Anc_G(\{u,v\} \cup S)$ by Lemma~\ref{lem:mixed_graph_represents_CDMG}.
  All colliders on some $\pi_{ij}$ are in $\Anc_G(\{i,j\}\cup S)$.
  Hence, they are in $\Anc_G(\{u,v\}\cup S)$.
  So, all colliders on $\pi$ are in $\Anc_G(\{u,v\}\cup S)$.
  All non-endpoint non-colliders on some $\pi_{ij}$ are unblockable, and therefore all non-endpoint non-colliders on $\pi$ are unblockable.
  Therefore, $\pi$ is a $\sigma$-inducing walk given $S$ in $G$.
  So there must also be a $\sigma$-inducing path given $S$ in $G$ between $u$ and $v$.
  Because $H$ represents $G$ given $S$, we conclude that $u,v$ must be adjacent in $H$.
\end{proof}
%\begin{Lem}\label{lem:into_undirected_edge}
%  Let $H = \lt J,V,E\rt$ be a mixed graph with input nodes $J$, output nodes $V$ and edges $E$ of the types $\{\tuh,\hut,\huo,\huh,\ouo,\ouh,\tuo,\out,\tut\}$.
%  Let $G$ be a CDMG with input nodes $J$ and output nodes $V^+ = V \dcup S$.
%  If $H$ represents $G$ given $S$, then no unshielded triple of the form $i \suh j \tut k$ can occur in $H$.
%\end{Lem}
%\begin{proof}
%  By contradiction.
%  Since $j \in \Anc_G(\{k\} \cup S)$ but $j \notin \Anc_G(\{i\} \cup S)$, we must have $j \in \Anc_G(k)$.
%  Also, $k \in \Anc_G(\{j\} \cup S)$.
%  Since $j \in \Anc_G(k)$ and $j \notin \Anc_G(S)$, we must have $k \in \Anc_G(j)$.
%  But then $j \in \Sc_G(k)$, which is not possible according to Lemma~\ref{lem:inducing_path_into_scc}.
%\end{proof}

The following result shows that every CDMG can be represented by a $\sigma$-PAG given some set of selection nodes.
\begin{Prp}\label{prp:from_cdmg_to_sigmapag}
  Let $G$ be a CDMG with input nodes $J$ and output nodes $V^+ = V \dcup S$. 
  There exists a $\sigma$-PAG, denoted $\sigmaPAG(G \given S)$, that represents $G$ given $S$.
\end{Prp}
\begin{proof}
  We will construct a mixed graph $H$ with input nodes $J$ and output nodes $V$ as follows.
  Let two nodes $i,j \in J\cup V$ be adjacent in $H$ if and only if $i \ne j$, $\{i,j\} \not\subseteq J$ and there is a $\sigma$-inducing path given $S$ between $i,j$ in $G$.
  In that case, orient the edge between $i$ and $j$ in $H$ as follows:
    $$\begin{cases}
      \text{$i \tut j$} & \text{if $i \in \Anc_G(\{j\} \cup S)$ and $j \in \Anc_G(\{i\} \cup S)$}, \\
      \text{$i \tuh  j$} & \text{if $i \in \Anc_G(\{j\} \cup S)$ and $j \notin \Anc_G(\{i\} \cup S)$}, \\
      \text{$i \hut  j$} & \text{if $i \notin \Anc_G(\{j\} \cup S)$ and $j \in \Anc_G(\{i\} \cup S)$}, \\
      \text{$i \huh j$} & \text{if $i \not\in \Anc_G(\{j\} \cup S)$ and $j \not\in \Anc_G(\{i\} \cup S)$.}
    \end{cases}$$
%    \begin{compactitem}
%      \item $u \to v \in \tilde{\C{E}}$ if and only if there is a $\sigma$-inducing path between $u$ and $v$ in $\C{G}$ and $u \in \ansub{\C{G}}{v}$,
%      \item $u \ot v \in \tilde{\C{E}}$ if and only if there is a $\sigma$-inducing path between $u$ and $v$ in $\C{G}$ and $v \in \ansub{\C{G}}{u}$,
%      \item $u \oto v \in \tilde{\C{F}}$ if and only if there is a $\sigma$-inducing path between $u$ and $v$ in $\C{G}$ and $u \not\in \ansub{\C{G}}{v}$ and $v \not\in \ansub{\C{G}}{u}$.
%    \end{compactitem}
%This construction preserves the (non-)ancestral relations as well as the $d$-separations/connections.
  It is obvious by construction that $H$ represents $G$ given $S$.
  It is a valid $\sigma$-PAG by Proposition~\ref{prp:sigmapag_represents_cdmg_is_valid}.
\end{proof}
The $\sigma$-PAG constructed in this way is a \emph{maximal $\sigma$-ancestral graph} ($\sigma$-MAG): it contains no circle edge marks and is therefore maximally informative about ancestral relations (that is, as informative as a $\sigma$-PAG can be).

\begin{Prp}\label{prp:inducing_edge_into_scc}
  Let $G$ be a CDMG with input nodes $J$ and output nodes $V^+ = V \dcup S$ and $H$ a $\sigma$-PAG that represents $G$ given $S$.
%If $i \huh j \sus k$ in $H$ and $j \in \Sc_G(k)$, then $i \huh k$ in $H$.
%If $i \tuh j \sus k$ in $H$ and $j \in \Sc_G(k)$, then either $i \tuh k$ or $i \huh k$ in $H$.
  If $i \suh j \sus k$ in $H$ (with $i,j,k$ distinct nodes) and $j \in \Sc_G(k)$, then $i \sus k$ in $H$ and $k \notin \Anc_G(\{i\} \cup S)$.
  If, additionally, the edge in $H$ between $i$ and $j$ is $i \huh j$, then also $i \notin \Anc_G(\{k\} \cup S)$.
\end{Prp}
\begin{proof}
Since $\Sc_G(k) \supseteq \{j,k\}$, $k$ cannot be in $J$.

By Lemma~\ref{lem:mixed_graph_represents_CDMG}, $i \suh j$ in $H$ implies the existence of a $\sigma$-inducing walk between $i$ and $j$ given $S$ in $G$ that is into $j$.
This can be extended by concatenation with a directed path from $j$ to $k$ into a $\sigma$-inducing walk between $i$ and $k$ given $S$ in $G$ that is into $k$.
Hence, there must be an edge $i \sus k$ in $H$.

%If it were of the form $i \sut k$, then $k \in \Anc_G(\{i\} \cup S)$.
If $k \in \Anc_G(\{i\} \cup S)$, then also $j \in \Anc_G(\{i\} \cup S)$ because $j \in \Anc_G(k)$, a contradiction.
%Hence the edge must be of the form $i \suh k$.

If $i \huh j$ in $H$, then $i \notin \Anc_G(\{j\} \cup S)$.
%If the edge between $i$ and $k$ were of the form $i \tuh k$, then $i \in \Anc_G(\{k\} \cup S)$.
If $i \in \Anc_G(\{k\} \cup S)$, then $i \in \Anc_G(k)$ because $i \notin \Anc_G(S)$; hence also $i \in \Anc_G(k) = \Anc_G(j)$, a contradiction.
%Hence if $i \huh j$ in $H$, the edge between $i$ and $k$ must be of the form $i \huh k$.
\end{proof}

The following important result states that the existence of a definitely open path in a $\sigma$-PAG representing a CDMG may imply the existence of a corresponding $\sigma$-open path in the CDMG.
\begin{Prp}\label{prp:sigmapag_open_walk_implies_cdmg_open_walk}
  Let $G$ be a CDMG with input nodes $J$ and output nodes $V^+ = V \dcup S$ and $H$ a $\sigma$-PAG that represents $G$ given $S$.
  Let, for $n \ge 2$,
    \[\pi = q_0 \sus q_1 \suh q_2 \huh \dots \huh q_{n-1} \hus q_n\] 
  be a definitely $Z$-open path in $H$, for $Z \subseteq (J \cup V) \sm \{q_0,q_n\}$,
  such that $q_2, \dots, q_{n-1}$ are definite colliders, and $q_{1}$ is either a definite non-collider or a definite collider.
  Then there exists a $(Z \cup S)$-$\sigma$-open path in $G$ between $q_0$ and $q_n$.
\end{Prp}
\begin{proof}
  The nodes $q_0, q_1, \dots, q_n$ in $J \cup V$ must all be distinct.
  For $i=1,\dots,n$, let $\mu_i$ be a $\sigma$-inducing path in $G$ between $q_{i-1}$ and $q_i$ that is into $q_{i-1}$ if the edge $q_i \hus q_{i-1}$ on $\pi$, and into $q_{i}$ if the edge $q_{i-1} \suh q_i$ on $\pi$ (see Lemma~\ref{lem:mixed_graph_represents_CDMG}).
  These can be concatenated into a walk $\mu = (\mu_1,\dots,\mu_n)$ in $G$ between $q_0$ and $q_n$:
  \[\overunderbraces{&\br{2}{\mu_1}&&&&\br{2}{\mu_{n}}}{&q_0 \sus \cdots \sus & q_1 & \sus \dots \suh q_2 & \hus \cdots \cdots \suh & q_{n-2} \hus \dots \suh & q_{n-1} & \hus \cdots \sus q_n}{&&\br{2}{\mu_2}&&\br{2}{\mu_{n-1}}}\]
  Let $Z \subseteq (J \cup V) \sm \{q_0,q_n\}$.
  Since $\pi$ is assumed to be definitely $Z$-open, $q_2,\dots,q_{n-1}$ are all in $\Anc_H(Z)$.
  We must have that $q_1 \in \Anc_G(\{q_0,q_2\} \cup Z \cup S)$, as can be seen by considering the two mutually exclusive cases:
  \begin{itemize}
    \item $q_1$ is a definite collider on $\pi$. Then $q_1 \in \Anc_H(Z)$. By Lemma~\ref{lem:mixed_graph_represents_CDMG}, $q_1 \in \Anc_G(Z \cup S)$.
    \item $q_1$ is a definite non-collider on $\pi$. Then $q_1 \notin Z$, but $q_1 \in \Anc_G(\{q_0,q_2\} \cup S)$ because either $q_0 \sut q_1$ or $q_1 \tus q_2$ must be on $\pi$.
  \end{itemize}
  In any case, $q_1 \in \Anc_G(\{q_0,q_2\} \cup Z \cup S)$.
  
  Consider all walks in $G$ between $q_0$ and $q_n$ with the property that 
  \begin{enumerate}
    \item all colliders on it are in $\Anc_{G}(\{q_0,q_n\} \cup Z \cup S)$, and 
    \item each non-endpoint non-collider on it is not in $\{q_0,q_n\} \cup Z \cup S$ or is unblockable.
  \end{enumerate}
  Such walks exist, since the concatenation $\mu = (\mu_1, \dots, \mu_n)$ is one, as we will now show. 

  To show that the first property holds for $\mu$, note that $q_2,\dots,q_{n-1}$ are in $\Anc_H(Z)$ and hence, by Lemma~\ref{lem:mixed_graph_represents_CDMG}, $q_2,\dots,q_{n-1} \in \Anc_G(Z \cup S)$, a subset of $\Anc_G(\{q_0,q_n\} \cup Z \cup S)$.
  We already saw that $q_1 \in \Anc_G(\{q_0,q_n\} \cup Z \cup S)$, which holds in particular if $q_1$ is a collider on $\mu$.
  Every internal collider on some $\mu_i$ is in $\Anc_G(\{q_{i-1},q_i\} \cup S)$, and hence also these `internal' colliders on $\mu$ are in $\Anc_G(\{q_0,q_n\} \cup Z \cup S)$.

  For the second property, note that all non-endpoint non-colliders on some $\mu_i$ are unblockable by assumption.
  $q_2,\dots,q_{n-1}$ cannot be non-colliders on $\mu$.
  If $q_1$ is a non-collider on $\mu$, then it must be a definite non-collider on $\pi$ and hence $q_1 \notin Z$ (and by assumption, $q_1 \notin S)$.
  Hence, the second property also holds for $\mu$.

  Let $\nu$ be a walk satisfying both properties above with a minimal number of colliders.
  We show that all colliders on $\nu$ must be in $\Anc_{G}(Z \cup S)$. 

  Suppose on the contrary the existence of a collider $k$ on $\nu$ that is not in $\Anc_G(Z \cup S)$. 
  It must then be in $\Anc_G(\{q_0,q_n\})$, by assumption.
  Then there exists a directed path in $G$ from $k$ to $q_0$ that does not pass through $q_n$, or there exists a directed path from $k$ to $q_n$ that does not pass through $q_0$.
  Without loss of generality, assume the former:
  there is a directed path in $G$ from $k$ to $q_0$ in $G$ that does not pass through $Z \cup S \cup \{q_n\}$.
  Let $\pi$ be a shortest path of that type.
  The directed path $\pi$ from $k$ to $q_0$ can be concatenated with the subwalk of $\nu$ between $k$ and $q_n$ (which is into $k$) into a walk between $q_0$ and $q_n$.
  This walk has the property, but has fewer colliders than $\nu$: a contradiction.

%  Therefore, $\nu$ is $\sigma$-open given $Z$ and all occurrences of $q_0$ and $q_n$ as non-endpoint non-collider are unblockable.
  Therefore, $\nu$ is a $(Z \cup S)$-$\sigma$-open walk in $G$ between $q_0$ and $q_n$.
  This means that there must also exist a $(Z \cup S)$-$\sigma$-open path in $G$ between $q_0$ and $q_n$.
\end{proof}

%\Joris{Tom suggested: perhaps one can extend the above to definitely open paths in $H$ of any form by changing the proof strategy a bit.
%If one encounters an undirected edge on the open path in $H$, and the two nodes are in the same scc in $G$, then one may be able to skip 
%the undirected edge by the spurious edge that has to exist. If one encounters an undirected edge on the open path in $H$ that signifies
%selection bias, there is nothing to show. Promising! This seems to work. If we first remove all undirected edges that correspond with
%two nodes in the same scc but not subject to selection bias (by just skipping one of the $q$'s, we can extend the $\sigma$-inducing 
%path to traverse the scc and keep the same edge orientation?????? ah it is more complicated; we may also encounter two consecutive
%undirected edges but then we can either substitute it by a single one in case of an scc or it signifies selection bias...). 
%So replace each maximal subwalk $q_i\dots q_j$ consisting entirely of undirected edges where all $q_i,\dots,q_j$ are in the same scc 
%either by a single undirected edge $q_i \tut q_j$ (if that maximal subwalk is actually the entire walk) or by a single edge
%$q_{i-1} \suh q_j$ if $q_{i-1} \suh q_i \tut \dots \tut q_j$ (keeping the orientation at $q_{i-1}$) or 
%$q_{i} \hus q_{j+1}$ if $q_i \tut \dots \tut q_j \hus q_{j+1}$ (keeping the edge orientation at $q_{j+1}$).
%Hm, what about $q_0 \hut q_1 \tut q_2 \tut \dots \tut q_{n-1} \tuh q_n$?
%The best we can do is replace it by $q_0 \hut q_1 \tut q_{n-1} \tuh q_n$.
%This strategy works: we can always follow tails to the left OR to the right until we hit
%a collider $q$ or an end node $q$, or an undirected edge between two $q$'s signifying selection bias. In all cases, we have established
%that the starting $q$ is ancestor of an end $q$ or $S$ or $Z$. So this holds for all $q$'s. Therefore, it also holds for all internal
%colliders. The non-colliders are easier to deal with.
%
%Interestingly, TOm said this was related to their Markov equivalence characterization in terms of triples: apparently
%one can skip some nodes on $\pi$ in order to get an uncovered path!
%
%Wait. It may not even be true. Take $X_1 \tuh X_2 \tut X_3 \hut X4$ with $X_2 \in \Sc_G(X_3)$ (two incoming edges on a 2-cycle).
%Then the path is open in $H$, but there is no open path in $G$. So this fails.}

\begin{figure}\centering
  \begin{tikzpicture}
    \begin{scope}
      \node at (-2,1.5) {(a)};
      \node[ndout] (X1) at (-1,1.5) {$i$};
      \node[ndout] (X2) at (-1,0) {$j$};
      \node[ndout] (X3) at (1,1.5) {$k$};
      \node[ndout] (X4) at (1,0) {$l$};
      \draw[tuh] (X1) edge (X2);
      \draw[tuh] (X3) edge (X4);
      \draw[tuh,bend left] (X2) edge (X4);
      \draw[tuh,bend left] (X4) edge (X2);
    \end{scope}
    \begin{scope}[xshift=5.5cm]
      \node at (-2,1.5) {(b)};
      \node[ndout] (X1) at (-1,1.5) {$i$};
      \node[ndout] (X2) at (-1,0) {$j$};
      \node[ndout] (X3) at (1,1.5) {$k$};
      \node[ndout] (X4) at (1,0) {$l$};
      \draw[tuh] (X1) edge (X2);
      \draw[tuh] (X3) edge (X4);
      \draw[tut] (X2) edge (X4);
      \draw[tuh] (X1) edge (X4);
      \draw[tuh] (X3) edge (X2);
    \end{scope}
%    \begin{scope}[xshift=0cm,yshift=-3cm]
%      \node at (-2,1.5) {(c)};
%      \node[ndout] (X1) at (-1,1.5) {$i$};
%      \node[ndout] (X2) at (-1,0) {$j$};
%      \node[ndout] (X3) at (1,1.5) {$k$};
%      \node[ndout] (X4) at (1,0) {$l$};
%      \node[ndout] (X5) at (0,-1) {$s$};
%      \draw[tuh] (X1) edge (X2);
%      \draw[tuh] (X3) edge (X4);
%      \draw[tuh] (X2) edge (X5);
%      \draw[tuh] (X4) edge (X5);
%      \draw[tuh] (X1) edge (X4);
%      \draw[tuh] (X3) edge (X2);
%    \end{scope}
%    \begin{scope}[xshift=5.5cm,yshift=-3cm]
%      \node at (-2,1.5) {(d)};
%      \node[ndout] (X1) at (-1,1.5) {$i$};
%      \node[ndout] (X2) at (-1,0) {$j$};
%      \node[ndout] (X3) at (1,1.5) {$k$};
%      \node[ndout] (X4) at (1,0) {$l$};
%      \draw[tut] (X1) edge (X2);
%      \draw[tut] (X3) edge (X4);
%      \draw[tut] (X2) edge (X4);
%      \draw[tut] (X1) edge (X4);
%      \draw[tut] (X3) edge (X2);
%    \end{scope}
  \end{tikzpicture}
  \caption{Example that shows that not all definitely open paths in a $\sigma$-PAG that represents a CDMG imply the existence of a corresponding $\sigma$-open path in the CDMG.
  (a) CDMG in which $i \Perp^\sigma k$, represented by the $\sigma$-PAG in (b) in which a definitely open path $i \tuh j \tut l \hut k$ exists.
  %(c) A CDMG that is represented by the $\sigma$-PAG in (d) given $\{s\}$ where $i \nPerp^\sigma k \given s$.
  \label{fig:sigmapags_are_no_ci_model}}
\end{figure}


\subsection{Unshielded triples}

One of the key steps in the FCI algorithm is the orientation of ``unshielded triples''.
The following proposition will later be used to ``orient'' the edges in unshielded triples in a $\sigma$-PAG representing a CDMG.
\begin{Prp}\label{prp:orienting_unshielded_triples}
  Let $G$ be a CDMG with input nodes $J$ and output nodes $V^+ = V \dcup S$ and $H$ a $\sigma$-PAG that represents $G$ given $S$.
  If $\lt i, j, k \rt$ form an unshielded triple in $H$ with $i \in V$ (and $j,k \in V \cup J$), then either 
  \begin{enumerate}[label=(\roman*)]
    \item $j \notin Z$ for each $Z \subseteq (V \cup J) \sm \{i,k\}$ such that $i \Perp^\sigma_G k \given Z \cup S$, and $j \notin \Anc_G(\{i,k\} \cup S)$, or
    \item $j \in Z$ for each $Z \subseteq (V \cup J) \sm \{i,k\}$ such that $i \Perp^\sigma_G k \given Z \cup S$, and $j \in \Anc_G(\{i,k\} \cup S)$.
  \end{enumerate}
\end{Prp}
\begin{proof}
  Case (i): $j \notin \Anc_G(\{i,k\} \cup S)$. 
  By orienting the edge marks at $j$ on the path $i \sus j \sus k$ in $H$ into $i \suh j \hus k$ (if not already oriented that way), we obtain a $\sigma$-PAG $\tilde{H}$ that represents $G$.
  Let $Z \subseteq (V \cup J) \sm \{i,k\}$ such that $i \Perp^\sigma_G k \given Z \cup S$. 
  If $j \in Z$, the path $i \suh j \hus k$ would be definitely $Z$-open in $\tilde{H}$.
  By Proposition~\ref{prp:sigmapag_open_walk_implies_cdmg_open_walk}, this implies the existence of a $(Z \cup S)$-$\sigma$-open walk in $G$ between $i$ and $k$, a contradiction.
  Hence $j \notin Z$.

  Case (ii): $j \in \Anc_G(\{i,k\} \cup S)$.
  By orienting the edge marks at $j$ on the path $i \sus j \sus k$ in $H$ into $i \sut j \sus k$ (if $j \in \Anc_G(\{i\} \cup S)$) or $i \sus j \tus k$ (if $j \in \Anc_G(\{k\} \cup S)$) or $i \sut j \tus k$ (if $j \in \Anc_G(\{i\} \cup S) \cap \Anc_G(\{k\} \cup S)$), we obtain a $\sigma$-PAG $\tilde{H}$ that represents $G$.
  Let $Z \subseteq (V \cup J) \sm \{i,k\}$ such that $i \Perp^\sigma_G k \given Z \cup S$. 
  If $j \notin Z$, the path $i \sus j \sus k$ is definitely $Z$-open in $\tilde{H}$.
  By Proposition~\ref{prp:sigmapag_open_walk_implies_cdmg_open_walk}, this implies the existence of a $(Z \cup S)$-$\sigma$-open walk in $G$ between $i$ and $k$, a contradiction.
  Hence $j \in Z$.
\end{proof}
Note that in the first case, we can orient the edges as $i \suh j \hus k$ (if they were not already oriented in this way) to obtain a $\sigma$-PAG $\tilde{H}$ that also represents $G$. 
In the second case, we cannot orient the edges, since we don't know whether $j \in \Anc_G(\{i\}\cup S)$ or $j \in \Anc_G(\{k\} \cup S)$ or both.

\subsection{Discriminating paths}

Another step in the FCI algorithm is related to the notion of ``discriminating paths''.
This can be considered as an extension of the notion of unshielded triple.
\begin{Def}\label{def:discriminating_path}
  A path $\pi = \lt i, j, q_1, \dots, q_n, k \rt$ (with $n \ge 1$) in a mixed graph $H$ is a \emph{discriminating path for $j$} if:
  \begin{enumerate}[label=(\roman*)]
    \item $i$ is not adjacent to $k$ in $H$, and
    \item for $r=1,\dots,n$: $q_r$ is a definite collider on $\pi$ and a parent of $i$ in $H$.
  \end{enumerate}
\end{Def}
Figure~\ref{fig:discriminating_path} illustrates this notion.

\begin{figure}\centering
  \begin{tikzpicture}
    \begin{scope}
      \node[ndout] (X) at (9,0) {$k$};
      \node[ndint] at (9,0) {$k$};
      \node[ndout] (Q1) at (7.5,0) {$q_n$};
      \node[ndout] (Q2) at (6,0) {$q_{n-1}$};
      \node (dots) at (4.5,0) {$\cdots$};
      \node[ndout] (Qn) at (3,0) {$q_1$};
%      \node[ndout] (B) at (1.5,0.0) {$b$};
%      \node[ndout] (Y) at (1.5,2) {$y$};
      \node[ndout] (B) at (1.5,0) {$j$};
      \node[ndint]     at (1.5,0) {$j$};
      \node[ndout] (Y) at (0,0) {$i$};
      \draw[suh] (X) -- (Q1);
      \draw[huh] (Q1) -- (Q2);
      \draw[huh] (Q2) -- (dots);
      \draw[huh] (dots) -- (Qn);
      \draw[tuh,bend right,in=240] (Q1) edge (Y);
      \draw[tuh,bend right,in=225] (Q2) edge (Y);
      \draw[tuh,bend right,in=210] (Qn) edge (Y);
      \draw[suh]  (B) -- (Qn);
      \draw[sus] (B) -- (Y);
    \end{scope}
  \end{tikzpicture}
  \caption{Discriminating path $\lt i,j,q_1,\dots,q_n,k \rt$ for $j$ between $i$ and $k$.
  Only $j$ and $k$ could be an input node, all other nodes must be output nodes.\label{fig:discriminating_path}}
\end{figure}

\begin{Rem}\label{rem:discriminating_path}
  It is instructive to think about a discriminating path rather as a certain collection of paths:
  \begin{align*}
    &i \hut q_n \hus k \\
    &i \hut q_{n-1} \huh q_n \hus k \\
    &i \hut q_{n-2} \huh q_{n-1} \huh q_n \hus k \\
    &\vdots \\
    &i \hut q_1 \huh \dots \huh q_{n-1} \huh q_n \hus k  \\
    &i \sus j \suh q_1 \huh \dots \huh q_{n-1} \huh q_n \hus k
  \end{align*}
  with the additional requirement that $i$ and $k$ are not adjacent.
\end{Rem}

The following quintessential property of discriminating paths is analogous to that of unshielded triples. 
\begin{Prp}\label{prp:orienting_discriminating_paths}
  Let $G$ be a CDMG with input nodes $J$ and output nodes $V^+ = V \dcup S$ and $H$ a $\sigma$-PAG that represents $G$ given $S$.
  If $\lt i, j, q_1, \dots, q_n, k \rt$ is a discriminating path in $H$ for $j$ between $i$ and $k$, then
  either 
  \begin{enumerate}[label=(\roman*)]
    \item $j \notin Z$ for each $Z \subseteq (V \cup J) \sm \{i,k\}$ such that $i \Perp^\sigma_G k | Z \cup S$, and $j \notin \Anc_G(\{q_1,i\} \cup S)$, or
    \item $j \in Z$ for each $Z \subseteq (V \cup J) \sm \{i,k\}$ such that $i \Perp^\sigma_G k | Z \cup S$, and $j \in \Anc_G(\{i\} \cup S)$.
  \end{enumerate}
  In both cases, $i \notin \Anc_G(\{j\} \cup S)$.
\end{Prp}
\begin{proof}
  Since $q_1$ is a parent of $i$ in $H$, $q_1 \in \Anc_G(\{i\} \cup S)$.
  Because $q_1$ is a definite collider in $H$, $q_1 \notin \Anc_G(S)$.
  Hence $q_1 \in \Anc_G(i) \sm \{i\}$, which means that $i \in V$.
  Also, this implies that $i \notin \Anc_G(j)$; otherwise, $q_1 \in \Anc_G(j)$ which contradicts $j \suh q_1$ in $H$.
  Hence $i \notin \Anc_G(\{j\} \cup S)$.

  We first show that if $Z \subseteq (V \cup J) \sm \{i,k\}$ such that $i \Perp^\sigma_G k \given Z \cup S$, then $\{q_1, \dots, q_n\} \in Z$.
  This can be seen from the various subpaths in Remark~\ref{rem:discriminating_path}.
  From the first path: if $q_n \notin Z$ then this path would be definitely open in $H$, and hence (by Proposition~\ref{prp:sigmapag_open_walk_implies_cdmg_open_walk}) there must be a $(Z \cup S)$-$\sigma$-open walk in $G$ between $i$ and $k$, which would contradict $i \Perp^\sigma_G k \given Z \cup S$.
  Once we have shown that $\{q_n, q_{n-1}, \dots, q_{n-r+1}\} \in Z$ for $1 \le r < n$, we see from the $r+1$'th path that $q_{n-r} \in Z$ to avoid definitely opening up this path in $H$, and thereby avoiding the existence of a $(Z\cup S)$-$\sigma$-open path in $G$ between $i$ and $k$.

  Case (i): $j \notin \Anc_G(\{q_1,i\} \cup S)$.
  By orienting the edge marks at $j$ on the path $i \sus j \sus q_1$ in $H$ into $i \suh j \hus q_1$ (if not already oriented that way), we obtain a $\sigma$-PAG $\tilde{H}$ that represents $G$.
%  Case (i): $j$ is a collider on the discriminating path.
  Let $Z \subseteq (V \cup J) \sm \{i,k\}$ such that $i \Perp^\sigma_G k \given Z \cup S$. 
  If $j \in Z$, the discriminating path in $\tilde{H}$ is definitely $Z$-open.
  By Proposition~\ref{prp:sigmapag_open_walk_implies_cdmg_open_walk}, this implies the existence of a $(Z \cup S)$-$\sigma$-open walk in $G$ between $i$ and $k$.
  Contradiction. 
  Hence $j \notin Z$.

  Case (ii): $j \in \Anc_G(\{q_1,i\} \cup S)$. 
  Since $q_1 \in \Anc_G(i)$, this implies $j \in \Anc_G(\{i\} \cup S)$.
%  Because of the edge $q_1 \tuh k$ in $H$, we also know that $k \notin \Anc_G(\{q_1\} \cup S)$.
%  Hence $j \in \Anc_G(k)$.
  We already concluded above that $i \notin \Anc_G(\{j\} \cup S)$.
  By orienting the the edge $i \sus j$ in $H$ as $i \hut j$, we obtain a $\sigma$-PAG $\tilde{H}$ that represents $G$.
%  Case (ii): $j$ is a non-collider in $H$ on the path $i \sus j \sus k$.
  Let $Z \subseteq (V \cup J) \sm \{i,k\}$ such that $i \Perp^\sigma_G k \given Z \cup S$. 
  If $j \notin Z$, the discriminating path in $\tilde{H}$ is definitely $Z$-open.
  By Proposition~\ref{prp:sigmapag_open_walk_implies_cdmg_open_walk}, this implies the existence of a $(Z \cup S)$-$\sigma$-open walk in $G$ between $i$ and $k$.
  Contradiction. 
  Hence $j \in Z$.
\end{proof}
Note that in the first case, we can orient $i \huh j \huh q_1$ (as far as the edge marks were not already oriented in this way) to obtain a $\sigma$-PAG $\tilde{H}$ that also represents $G$. 
In the second case, we can orient $i \hut j$ (as far as the edge marks were not already oriented in this way) to obtain a $\sigma$-PAG $\tilde{H}$ that also represents $G$.

\subsection{Independence models and Markov equivalence}

We start with an abstract definition of an ``independence model'', where we extend the common definition to allow for input nodes.
\begin{Def}\label{def:independence_model}
  Given two disjoint sets $J,V$ (the `inputs' and `outputs', respectively), we call a subset of 
  $$\{ \lt A, B, C \rt : A, B, C \subseteq J \cup V, A \cap J = \emptyset, J \subseteq B \cup C  \}$$
  an \emph{independence model over $V \given J$}.
  For an element $\lt A, B, C \rt$ of an independence model, we also say that \emph{$C$ separates $A$ from $B$}.
\end{Def}

Independence models can be used to encode all (conditional) independences in a Markov kernel.
\begin{Def}\label{def:independence_model_Markov_kernel}
  For a Markov kernel $K(X_V|X_J) : \C{X}_J \dshto \C{X}_V$, we define its \emph{independence model} to be
  $$\IM(K(X_V | X_J)) := \{ \lt A, B, C \rt : A, B, C \subseteq J \cup V, A \cap J = \emptyset, J \subseteq B \cup C : X_A \Indep_{K(X_V | X_J)} X_B \mid X_C \},$$
  i.e., the set of all conditional independences (of restricted form) that $K(X_V|X_J)$ satisfies.
\end{Def}

Independence models can also encode all separations in a graph.
\begin{Def}\label{def:independence_model_graph}
  For a CDMG $G$ with input nodes $J$ and output nodes $V$, define its \emph{$\sigma$-independence model} to be
  $$\IM_\sigma(G) := \{ \lt A, B, C \rt : A, B, C \subseteq J \cup V, A \cap J = \emptyset, J \subseteq B \cup C : A \Perp_G^\sigma B \mid C \},$$
  i.e., the set of all $\sigma$-separations (of restricted form) entailed by the graph. 
  Define its \emph{$d$-independence model} to be
  $$\IM_d(G) := \{ \lt A, B, C \rt : A, B, C \subseteq J \cup V, A \cap J = \emptyset, J \subseteq B \cup C : A \Perp_G^d B \mid C \},$$
  i.e., the set of all $d$-separations (of restricted form) entailed by the graph.
\end{Def}
Both $\IM_d(G)$ and $\IM_\sigma(G)$ are independence models over $V \given J$ (with $J$ the input nodes of $G$ and $V$ the output nodes of $G$).
For CADMGs, $\sigma$-separation is equivalent to $d$-separation, and hence, if $G$ is acyclic, then $\IM_d(G) = \IM_\sigma(G)$.

When conditioning on a set of selection variables, we can also define conditional independence models from graphs as follows.
\begin{Def}\label{def:conditional_independence_model_graph}
  For a CDMG $G$ with input nodes $J$ and output nodes $V^+ = V \dcup S$, define its \emph{$\sigma$-independence model given $S$} to be
  $$\IM_\sigma(G \given S) := \{ \lt A, B, C \rt : A, B, C \subseteq J \cup V, A \cap J = \emptyset, J \subseteq B \cup C : A \Perp_G^\sigma B \mid C \cup S \},$$
  i.e., the set of all $\sigma$-separations (of restricted form) involving nodes not in $S$, when also conditioning on $S$. 
  Define its \emph{$d$-independence model given $S$} to be
  $$\IM_d(G \given S) := \{ \lt A, B, C \rt : A, B, C \subseteq J \cup V, A \cap J = \emptyset, J \subseteq B \cup C : A \Perp_G^d B \mid C \cup S \},$$
  i.e., the set of all $d$-separations (of restricted form) involving nodes not in $S$, when also conditioning on $S$.
\end{Def}
Hence, also $\IM_\sigma(G \given S)$ and $\IM_d(G \given S)$ are independence models over $V \given J$ (with $J$ the input nodes of $G$ and $V$ the output nodes of $G$ except for those in $S$).
For CADMGs, $\sigma$-separation is equivalent to $d$-separation, and hence, if $G$ is acyclic, then $\IM_d(G \given S) = \IM_\sigma(G \given S)$.

%To connect with the literature, we also provide the following definition in case no exogenous input nodes are present:
%\begin{Def}\label{def:independence_model}
%  For a DMG $G$ with nodes $V$, define its \emph{$d$-independence model} to be
%  $$\IM_d(G) := \{ \lt A, B, C \rt : A, B, C \subseteq V, A \Perp_G^d B \mid C \},$$
%  i.e., the set of all $d$-separations entailed by the graph. Define its
%  \emph{$\sigma$-independence model} to be
%  $$\IM_\sigma(G) := \{ \lt A, B, C \rt : A, B, C \subseteq V, A \Perp_G^\sigma B \mid C \},$$
%  i.e., the set of all $\sigma$-separations entailed by the graph. 
%\end{Def}
%In general, given an index set $V$, we call a subset of 
%$\{ \lt A, B, C \rt : A, B, C \subseteq V \}$
%an \emph{independence model over $V$}.
%If $\lt A, B, C \rt$ in an independence model, we also say that $C$ separates $A$ from $B$.
%Both $\IM_d(G)$ and $\IM_\sigma(G)$ are independence models over $V$, the nodes of $G$.
%For ADMGs, $\sigma$-separation is equivalent to $d$-separation, and hence, if $G$ is acyclic, then
%$\IM_d(G) = \IM_\sigma(G)$.

%The input to the original FCI algorithm (which assumes $J = \emptyset$) consists of an independence model over $V$, and its output will consist of a mixed graph with nodes $V$.
%If the input of FCI is the independence model of a DMG, then the output of FCI will be a $\sigma$-PAG that represents the ``Markov equivalence class'' of the DMG.
%\begin{Def}\label{def:Markov_equivalent}
%  We call two DMGs $G_1$ and $G_2$ \emph{$\sigma$-Markov equivalent} if $\IM_\sigma(G_1) = \IM_\sigma(G_2)$,
%  and \emph{$d$-Markov equivalent} if $\IM_d(G_1) = \IM_d(G_2)$.
%\end{Def}

For notational convenience, we define symmetrized versions of independence models.
\begin{Def}\label{def:symmetrized_independence_model}
  Given two disjoint sets $J,V$ (the `inputs' and `outputs', respectively), 
  and an independence model $I$ over $V \given J$, we define the \emph{symmetrized} independence model as
  $$I^+ := I \cup \{ \lt B, A, C \rt : \lt A, B, C \rt \in I \}. $$
\end{Def}

The input to the extended FCI algorithm will consist of an independence model over $V \given J$, and its output will consist of a mixed graph with input nodes $J$ and output nodes $V$.
One of the nice properties of the extended FCI algorithm is that if the input of FCI is the independence model of a CDMG given $S$, then the output of FCI will be a $\sigma$-PAG that represents the ``Markov equivalence class'' of the CDMG (at least if $S = \emptyset$ or $J = \emptyset$).
\begin{Def}\label{def:Markov_equivalent}
  Let $G_1, G_2$ be two CDMGs with input nodes $J$ and output nodes $V_1^+ = V \dcup S_1$, $V_2^+ = V \dcup S_2$, respectively.
  We call $G_1$ and $G_2$ \emph{$\sigma$-Markov equivalent w.r.t.\ $V \given J$} if $\IM_\sigma(G_1 \given S_1) = \IM_\sigma(G_2 \given S_2)$, and \emph{$d$-Markov equivalent w.r.t.\ $V \given J$} if $\IM_d(G_1 \given S_1) = \IM_d(G_2 \given S_2)$.
  %In these cases, we write $G_1 \underset{V \given J}{\overset{\sigma}{\sim}} G_2$, respectively $G_1 \underset{V \given J}{\overset{d}{\sim}} G_2$.
\end{Def}

When exploiting conditional independences for causal discovery, one typically has to make some kind of faithfulness assumption.
The faithfulness assumption that we will make for the extended FCI algorithm will be that the (restricted) conditional independences in the conditioned Markov kernel are identical to the (restricted) $\sigma$-separations in the graph, given $S$:
$$\IM(P_{M}(X_V \mid X_S \in \xi_S, \doit(X_J))) = \IM_\sigma(G_{V^+|J}(M) \given S).$$

\subsection{Skeleton search}

The FCI algorithm consists of two phases, the skeleton search phase, which is followed by the orientation phase.
In this subsection, we describe the skeleton search phase.

\begin{Def}\label{def:skeleton}
Given a $\sigma$-PAG $H = \lt J, V, E \rt$, its \emph{skeleton} is the mixed graph 
$\skel(H) := \lt J, V, F \rt$ with the same nodes, and with a bicircle edge $i \ouo j$ in $F$
if and only if $i \sus j$ in $E$ (i.e., if $i$ and $j$ are adjacent in $H$).
\end{Def}
Hence, the only edge type occurring in the skeleton is the bicircle edge.
The skeleton has no edge between any pair of input nodes.
We will later frequently refer to the unordered pairs of nodes that may be adjacent:
\begin{Def}
  For input nodes $J$ and output nodes $V$, define 
  $$\separable(V|J) :=  \{ \{i,j\} : i \in V, j \in J \cup V, i \ne j \}.$$
  %\{ \{i,j\} : i \in V, j \in V, i \ne j \} \cup \{ (i,j) : i \in V, j \in J \}.
\end{Def}

The aim of the skeleton search phase of the FCI algorithm is to construct the skeleton of the $\sigma$-PAG that represents a CDMG given $S$ from the $\sigma$-independence model given $S$ of the CDMG.
It does this by testing for each separable edge in the skeleton whether it can find any subset of nodes that separates the two nodes.
If it finds a separating set between an (unordered) pair of distinct nodes $\{i, j\}$, the set is stored as $\sepset(\{i,j\})$.
The orientation phase later makes use of these separating sets found in the skeleton phase.

A brute-force search over all possible subsets of $(J \cup V)\sm \{i,j\}$, as in Algorithm~\ref{alg:bf_skeleton}, would be a straightforward solution.
However, it is also computationally extremely expensive for all but the smallest cardinalities of $V$ and $J$, and statistically not very reliable in case the separations have to be tested with conditional independence tests on finite data.
\begin{algorithm}
  \begin{algorithmic}[1]
    \Input Input node set $J$; output node set $V$; independence model $I$ over $V \given J$
    \Output mixed graph $H$ with input nodes $J$ and output nodes $V$; separating sets $\sepset$
    \State $H \leftarrow \lt J, V, \emptyset \rt$
    \ForEach{$(i,j) \in \separable(V|J)$}
      \State add edge $i \ouo j$ to $H$
      \ForEach{$Z \subseteq (V \cup J) \sm \{i,j\}$}
      \If{$\lt i, j, Z \rt \in I^+$}\Comment{found a separating set}
        \State delete edge $i \ouo j$ from $H$
        \State $\sepset(\{i,j\}) \leftarrow Z$
        \Break
      \EndIf
      \EndFor
    \EndFor
  \end{algorithmic}
  \caption{Brute-force skeleton algorithm.\label{alg:bf_skeleton}}
\end{algorithm}

To get some inspiration on how to address this, we will first describe the skeleton search phase of the PC algorithm, an ancestor of the FCI algorithm designed for DAGs \cite{SGS2000}.
The PC skeleton phase (Algorithm~\ref{alg:pc_skeleton}) searches for separating sets of increasing cardinality. 
As candidates for a separating set between nodes $i,j \in H$, it considers all subsets of $\Adj_H(i)$ and all subsets of $\Adj_H(j)$.
The rationale is that if $G$ is a DAG, then either $\Pa_G(i) \subseteq \Adj_H(i)$ or $\Pa_G(j) \subseteq \Adj_H(j)$ $d$-separates $i$ from $j$.

\begin{algorithm}
  \begin{algorithmic}[1]
  \Input Output node set $V$; independence model $I$ over $V \given \emptyset$
  \Output mixed graph $H$ with output nodes $V$; separating sets $\sepset$
  \State initialize $H$ as a complete graph with only bicircle ($\ouo$) edges
  \State $n \leftarrow 0$
    \Repeat
    \Repeat
    \State select $i,j \in V$ with $i \ouo j$ in $H$ and $\#(\Adj_H(i)\sm\{j\}) \ge n$
    \State select a subset $Z \subseteq \Adj_H(i)\sm\{j\}$ of cardinality $n$
    \If{$(i,j,Z) \in I^+$}
      \State delete edge $i \ouo j$ from $H$
      \State $\sepset(\{i,j\}) \leftarrow Z$
    \EndIf
    \Until{no more such tuples $(i,j,Z)$ can be selected}
    \State $n \leftarrow n + 1$
    \Until{for all $i,j \in V$ with $i \ouo j$ in $H$, $\#(\Adj_H(i)\sm\{j\}) < n$}
  \end{algorithmic}
  \caption{Original PC skeleton algorithm.\label{alg:pc_skeleton}}
\end{algorithm}

We can easily extend this to deal with input nodes as well, see Algorithm~\ref{alg:pc_skeleton_extended}.
We may restrict ourselves to search for separating sets that contain all nodes in $J$ (except $j$ itself, if $j \in J$).
We therefore only need to consider subsets of neighbours of $i$ that are not in $J$, in increasing cardinality.

\begin{algorithm}
  \begin{algorithmic}[1]
  \Input Input node set $J$; output node set $V$; independence model $I$ over $V \given J$
  \Output mixed graph $H$ with input nodes $J$ and output nodes $V$; separating sets $\sepset$
  \State $H \leftarrow \lt J, V, \emptyset \rt$
  \ForEach{$(i,j) \in \separable(V|J)$}
    \State add edge $i \ouo j$ to $H$
  \EndFor
  \State $n \leftarrow 0$
    \Repeat
    \Repeat
    \State select $i\in V, j \in V \cup J$ with $i \ouo j$ in $H$ and $\#(\Adj_H(i)\sm(J\cup\{j\})) \ge n$
    \State select a subset $Z \subseteq \Adj_H(i)\sm(J\cup\{j\})$ of cardinality $n$
    \If{$(i,j,Z \cup (J\sm\{j\})) \in I^+$} 
      \State delete edge $i \ouo j$ from $H$
      \State $\sepset(\{i,j\}) \leftarrow Z \cup (J \sm \{j\})$
    \EndIf
    \Until{no more such tuples $(i,j,Z)$ can be selected}
    \State $n \leftarrow n + 1$
    \Until{for all $i\in V, j \in V \cup J$ with $i \ouo j$ in $H$, $\#(\Adj_H(i)\sm(J\cup\{j\})) < n$}
  \end{algorithmic}
  \caption{Extended PC skeleton algorithm $\mathtt{PCskeleton}(J,V,I)$.\label{alg:pc_skeleton_extended}}
\end{algorithm}

The underlying idea may no longer hold if $G$ is a CADMG or a CDMG, or in case of selection bias.
The original proposal (which motivated the somewhat optimistic adjective ``Fast'' in the name of the FCI algorithm \cite{SMR1995}) replaces the subsets of adjacent nodes by so-called ``Possible-D-Sep'' sets.\footnote{An alternative search strategy was proposed that can be considerably faster in practice, for which it can be shown that the corresponding FCI+ algorithm is of polynomial-time complexity in the number of nodes, as long as the degree of the DPAG is bounded \cite{ClaassenMooijHeskes_UAI_13}.}
Before we can define these, we need some definitions and theory.

\begin{Def}\label{def:SEP}
  Let $G$ be a CDMG with input nodes $J$ and output nodes $V^+ = V \dcup S$.
  For distinct nodes $i,j \in V \dcup J$,
  define $\SEP_G(i,j)$ as the set of nodes $k \in V$ such that
  $k \ne i$ and there is a walk $\pi$ in $G$ between $i$ and $k$ such that every node
  on $\pi$ is in $\Anc_G(\{i,j\}\cup S)$, and every non-endpoint non-collider on $\pi$ is unblockable.
\end{Def}
The name $\SEP_G(i,j)$ is motivated by the following property.
\begin{Prp}\label{prp:SEP}
  Let $G$ be a CDMG with input nodes $J$ and output nodes $V^+ = V \dcup S$.
  Let $i \in V$ and $j \in J \cup V$ be distinct nodes.
  If there exists no $\sigma$-inducing walk given $S$ between $i,j$ in $G$, then $\SEP_G(i,j) \cap \{i,j\} = \emptyset$ and $i \Perp^\sigma_G j \given \SEP_G(i,j) \cup (J \sm \{j\}) \cup S$.
\end{Prp}
\begin{proof}
  By definition, $\SEP_G(i,j) \subseteq \Anc_G(\{i,j\} \cup S)$ and $i \notin \SEP_G(i,j)$.
  If $j \in \SEP_G(i,j)$, the walk $\pi$ in Definition~\ref{def:SEP} between $i$ and $j$ must be $\sigma$-inducing given $S$.
  Since there exists no $\sigma$-inducing walk given $S$ between $i,j$ in $G$ by assumption, $j \notin \SEP_G(i,j)$.

  We prove the $\sigma$-separation by contradiction.
  Write $Z := \SEP_G(i,j) \cup (J \sm \{j\}) \cup S$.
  Suppose there exists a path in $G$ between $i$ and $\{j\} \cup J$ that is $\sigma$-open given $Z$.
  It cannot end in a node in $J \sm \{j\}$, and hence it must end in $j$.
  Let $\pi = v_0 \sus \dots \sus v_n$ with $v_0 = i$ and $v_n \in j$ be such a path consisting of a minimal number of nodes.
  Every node on $\pi$ must be in $\Anc_G(\{i,j\} \cup S) \cup J$, since by Lemma~\ref{lem:ancestral_relations_open_path}, each node on $\pi$ is in $\Anc_G((\{i,j\} \sm J)\cup Z)$, and $Z \subseteq \Anc_G(\{i,j\} \cup S) \cup J$.
  $\pi$ can only contain nodes in $J$ as endnodes (otherwise we could shorten the path).
  In other words, the only node on $\pi$ that could be in $J$ is $j$.

  Denote the subwalk of $\pi$ from $v_a$ to (and including) $v_b$ by $\pi(v_a,v_b)$, for $0 \le a \le b \le n$.
  We will show that for all $k = 1,\dots,n$, $\pi(v_0,v_k)$ has the property that every non-endpoint non-collider on it is unblockable.
  The property trivially holds for $k=1$. 
  Suppose it holds for $k < n$.
  Since $v_0,\dots,v_k$ are all in $\Anc_G(\{i,j\} \cup S)$, and all non-endpoint non-colliders on $\pi(v_0,v_k)$ are unblockable, we conclude that $v_k \in \SEP_G(i,j)$.
  If $v_{k}$ is a non-collider on $\pi(v_0,v_{k+1})$, it must be unblockable, because $\pi(v_0,v_{k+1})$ is $Z$-$\sigma$-open and $v_k \in \SEP_G(i,j) \subseteq Z$.
  So the property also holds for $k+1$.

  In particular, we can conclude that $j = v_n \in \SEP_G(i,j)$, a contradiction.
\end{proof}
In practice, one does not know the set $\SEP_G(i,j)$ if $G$ is unknown.
One can, however, easily obtain a `bound' on this set by identifying a superset.
\begin{Def}\label{def:possep}
  Let $H_0 = \lt J,V,E\rt$ be a mixed graph with input nodes $J$, output nodes $V$ and edges $E$ of the types $\{\tuh,\hut,\huo,\huh,\ouo,\ouh,\tuo,\out,\tut\}$.
  For $i \in V, j \in J \cup V$ distinct, we define $\PosSEP_{H_0}(i,j) \subseteq V$ to consist of those nodes
  $k \in V$ such that $k \not\in \{i,j\}$ and there is a path between $i$ and $k$ in $H_0$ such that for
  every subsequent triple $a \sus b \sus c$ on the path, either the triple is a definite collider in $H_0$,
  or a triangle in $H_0$.
\end{Def}
We can now show that:
\begin{Lem}\label{lem:possep}
  Let $G$ be a CDMG with input nodes $J$ and output nodes $V^+ = V \dcup S$.
  Let $H_0$ be the mixed graph constructed by lines 
  \ref{alg:fci_skeleton2:cpc_skeleton}--\ref{alg:fci_skeleton2:oriented_unshielded_colliders}
  of Algorithm~\ref{alg:fci_skeleton} when run on the independence model $\IM_\sigma(G \given S)$.
  Then, for $i \in V$ and $j \in J \cup V$ distinct nodes, if there is no $\sigma$-inducing path given $S$ in $G$ between $i$ and $j$, then
  $\SEP_G(i,j) \subseteq \PosSEP_{H_0}(i,j)$.
\end{Lem}
\begin{proof}
%  That is, first one runs the Conditional PC skeleton algorithm (Algorithm~\ref{alg:pc_skeleton}),
%  then one orients all unshielded colliders as definite colliders if the middle node is in the separating set of the two outer mnodes, and finally one orients the edges between input and output nodes.
  Note that the skeleton of $H_0$ is a supergraph of the skeleton of $\sigmaPAG(G \given S)$, i.e., every adjacency in $\sigmaPAG(G \given S)$ must also be an adjacency in $H_0$.
  Let $k \in \SEP_G(i,j)$.
  Since there is no $\sigma$-inducing path given $S$ in $G$ between $i$ and $j$, $\{i,j\} \cap \SEP_G(i,j) = \emptyset$ by Proposition~\ref{prp:SEP}, and hence $k \ne i$ and $k \ne j$.
  By definition, there is a path $\pi$ in $G$ between $i$ and $k$ with the property that every node on $\pi$ is in $\Anc_G(\{i,j\}\cup S) \cap V$, and every non-endpoint non-collider on $\pi$ is unblockable.
  There is a corresponding path $\pi'$ in $H_0$ between $i$ and $k$ that consists of the same sequence of nodes (but may have different edges between the nodes).

  Consider a subsequent triple $a \sus b \sus c$ on $\pi$.
  Suppose $b$ is a collider on $\pi$.
  In $H_0$, $a$ must also be adjacent to $b$, and $b$ to $c$.
  If $a$ and $c$ are not adjacent in $H_0$, then a separating set for $a$ and $c$ was found during the construction of $H_0$, or $a$ and $c$ are both in $J$.
  The latter would be a contradiction. 
  Therefore, it will have been oriented as a definite collider in $H_0$ according to Proposition~\ref{prp:orienting_unshielded_triples}.
  Otherwise (if $a$ and $c$ are adjacent in $H_0$), it forms a triangle in $H_0$.
  If $b$ is a non-collider on $\pi$, then it must be unblockable.
  But in that case, $a \sus b \sus c$ is a $\sigma$-inducing walk given $S$ between $a$ and $c$ in $G$.
  Hence, $(a,b,c)$ will form a triangle in $H_0$.
  Therefore, $k \in \PosSEP_{H_0}(i,j)$.
\end{proof}

\begin{algorithm}
  \begin{algorithmic}[1]
  \Input Input node set $J$; output node set $V$; independence model $I$ over $V \given J$
  \Output mixed graph $H$ with input nodes $J$ and output nodes $V$; separating sets $\sepset$
  \State $\lt H_0,\sepset \rt \leftarrow \mathtt{PCskeleton}(J,V,I)$ \label{alg:fci_skeleton2:cpc_skeleton}
%  \For{$j \ouo v$ in $H$ with $j \in J$, $v \in V$} \Comment{apply background knowledge \Joris{REDUNDANT?}}
%    \State orient the edge as $j \tuh v$
%  \EndFor
%  \State Orient unshielded triples in $H_0$:
%    \begin{description}
%      \item[$\C{R}0$] for each unshielded triple $\lt i,j,k \rt$ in $H_0$ with $j,k \in V$, orient it as $i \suh j \hus k$ if $j \notin \sepset(k,i)$
%    \end{description}\label{alg:fci_skeleton2:oriented_unshielded_colliders}
  \ForEach{unshielded triple $\lt i,j,k \rt$ in $H_0$ with $i,j \in V$} \Comment{orient colliders}
    \If{$j \notin \sepset(\{i,k\})$} % \sepset(k,i) if ordered!
      \State orient it as $i \suh j \hus k$
    \EndIf
  \EndFor\label{alg:fci_skeleton2:oriented_unshielded_colliders}
%  \For{$i \sus j \sus k$ in $H_0$ with $j,k \in V$}  \Comment{orient unshielded colliders}
%    \If{$i \notin \Adj_{H_0}(k)$}
%      \If{$j \notin \sepset(k,i)$}
%        \State orient it as $i \suh j \hus k$
%      \EndIf
%    \EndIf
%  \EndFor \label{alg:fci_skeleton2:oriented_unshielded_colliders}
  %find PossibleDSep sets
  \ForEach{$i \in V, j \in V \cup J$ with $i \ne j$} \Comment{construct Possible-D-Sep sets}
    \State calculate $\PosSEP_{H_0}(i,j)$
  \EndFor
  %remove orientations
    \State $H \leftarrow (J,V,\{i \ouo j : i \sus j \in H_0 \})$ \Comment{forget orientations} \label{alg:fci_skeleton2:construct_skeleton_start}
%  \For{$(i,j) \in \separable(V|J)$} \Comment{remove orientations}
%    \If{$i \sus j$ in $H_0$}
%      \State replace the edge by $i \ouo j$
%    \EndIf
%  \EndFor
  %find more separating sets
  \State $n \leftarrow 0$
    \Repeat \Comment{search for separating sets}
    \Repeat
    \State select $i \in V, j \in V \cup J$ with $i \ouo j$ in $H$ and $\#(\PosSEP_{H_0}(i,j)) \ge n$
    \State select a subset $Z \subseteq \PosSEP_{H_0}(i,j)$ of cardinality $n$
    \If{$(i,j,Z \cup (J\sm\{j\})) \in I$} 
      \State delete edge $i \ouo j$ from $H$
      \State $\sepset(\{i,j\}) \leftarrow Z \cup (J \sm \{j\})$
    \EndIf
    \Until{no more such tuples $(i,j,Z)$ can be selected}
    \State $n \leftarrow n + 1$
    \Until{for all $i \in V, j \in V \cup J$ with $i \ouo j$ in $H$, $\#(\PosSEP_{H_0}(i,j)) < n$}\label{alg:fci_skeleton2:construct_skeleton_end}
  \end{algorithmic}
  \caption{Extended FCI skeleton algorithm $\mathtt{FCIskeleton}(J,V,I)$.\label{alg:fci_skeleton}}
\end{algorithm}
%The FCI skeleton search uses the strategy of the PC skeleton algorithm, but uses $\PosSEP_{H_0}(i,j)$ instead of $\Adj_H(i)\sm\{j\}$ and check for each subset whether it separates $i$ from $j$ (when conditioning also on $J \sm \{j\}$). 
%Lemma~\ref{lem:possep} guarantees that such a separating set exists, if any separating set of $i$ and $j$ exists. 

So we do not need to search over all possible subsets of $(J \cup V) \sm \{i,j\}$ for a separating set between $i,j$, but only the subsets in $\PosSEP_{H_0}(i,j)$.
The skeleton phase of the extended FCI algorithm is described in Algorithm~\ref{alg:fci_skeleton}. 
It is an extension of the original FCI skeleton search \cite{SMR1995} to deal with input nodes. 
It first runs the extended PC skeleton phase (Algorithm~\ref{alg:pc_skeleton_extended}) and orients unshielded triples, obtaining a directed mixed graph $H_0$. 
Then it calculates the sets $\PosSEP_{H_0}(i,j)$ for all distinct pairs $(i,j)$ with $i \in V, j \in J \cup V$.
If $i \Perp^\sigma_G j \mid Z \cup S$ for some $Z \subseteq (J \cup V) \sm \{i,j\}$, then
$i \Perp^\sigma_G j \mid Z^* \cup S$ for some $Z^* \subseteq \PosSEP_{H_0}(i,j)$.
Since some of the oriented colliders may be incorrect in $H_0$ (because the PC skeleton phase may not have found all separating sets), it removes all orientations from $H_0$ and then continues with a more extensive search for separating sets, similar
to how the PC skeleton phase is done, but now using $\PosSEP_{H_0}(i,j)$ instead of $\Adj_{H_0}(i)\sm\{j\}$ to find candidate nodes for the separating set.

\begin{Thm}\label{thm:fci_skeleton_sound}
  The extended FCI skeleton algorithm (Algorithm~\ref{alg:fci_skeleton}) is \emph{sound}: if its input consists of the $\sigma$-independence model $\IM_\sigma(G \given S)$ given $S$ of a CDMG $G$, then its output will be $\skel(\sigmaPAG(G \given S))$. Furthermore, $i \Perp^\sigma j \given \sepset(\{i,j\})$ for all $i \in V, j \in V \cup J$ for node pairs $(i,j) \in \separable(V|J)$ that are non-adjacent in the skeleton.
\end{Thm}
\begin{proof}
  Let $G$ be a CDMG with input nodes $J$ and output nodes $V^+ = V \dcup S$ and $I = \IM_\sigma(G \given S)$ its $\sigma$-independence model given $S$.
  By Lemma~\ref{lem:possep}, the mixed graph $H_0$ constructed by lines 
  \ref{alg:fci_skeleton2:cpc_skeleton}--\ref{alg:fci_skeleton2:oriented_unshielded_colliders}
  has the property that for distinct $i \in V, j \in J \cup V$, if there is no $\sigma$-inducing path given $S$ in $G$ between $i$ and $j$, then $\SEP_G(i,j) \subseteq \PosSEP_{H_0}(i,j)$.
  Thus, if $i \in V$ and $j \in J \cup V$ are not adjacent in $\skel(\sigmaPAG(G \given S))$, then we are guaranteed that for some subset $Z \subseteq \PosSEP_{H_0}(i,j)$, we have that $i \Perp^\sigma_G j \given (Z \cup J) \sm \{j\} \cup S$.
%  We have to show that which means that it will compute $H = \skel(\sigmaPAG(G))$, and $\sepset$ will contain a separating set of $i$ and $j$ for every edge $i \ouo j$ absent in $H$ with $i \in V, j \in J \cup V, i \ne j$.
  The graph $H$ constructed in lines \ref{alg:fci_skeleton2:construct_skeleton_start}--\ref{alg:fci_skeleton2:construct_skeleton_end} will be a mixed graph with input nodes $J$ and output nodes $V$ that only has bicircle edges.
  It has an edge between any pair of distinct nodes $i,j \in J\cup V$ if and only if $\{i,j\} \not\ins J$ and there is no set $Z \subseteq (J \cup V) \sm \{i,j\}$ such that $i \Perp^\sigma_G j \given Z \cup S$.
  Furthermore, if there is no edge in $H$ between a pair of distinct nodes $i,j \in J \cup V$, then either $\{i,j\} \subseteq J$, or $i \Perp^\sigma_G j \given \sepset(\{i,j\}) \cup S$.
  By Proposition~\ref{prp:sigma-inducing_path_properties}, this implies that two distinct nodes $i,j \in J \cup V$ are adjacent in $H$ at this stage if and only if there is a $\sigma$-inducing walk given $S$ in $G$ between $i$ and $j$.
  Hence $H$ must be $\skel(\sigmaPAG(G \given S))$.
\end{proof}

%Another variant, the ``Really Fast Causal Inference'' algorithm has also been proposed, which adapts both the skeleton phase and the orientation rules, and sacrifices some completeness but remains sound \cite{Colombo++2012}.

\subsection{FCI Algorithm}

We are now ready to describe a causal inference algorithm that is an extension of the original
Fast Causal Inference (FCI) algorithm of \cite{SMR1995} to deal with input nodes and cycles.
It is presented as Algorithm~\ref{alg:fci}.
Its input is an independence model over $V \given J$, where $V$ and $J$ are index sets of output and input nodes, respectively. 
Its output is a mixed graph with input nodes $J$ and output nodes $V$.
It starts with a \emph{skeleton phase} (line~\ref{alg:fci:skeleton}, see Algorithm~\ref{alg:fci_skeleton}) that is aimed at deducing the adjacencies between the nodes, and to find sets that separate two separable node pairs.
Then, it runs various \emph{orientation rules} (lines~\ref{alg:fci:start_orient}--\ref{alg:fci:end_orient}) that iteratively orient circle edge marks into tails and arrowheads.
Note that by convention, the labeled nodes within each orientation rule are assumed to be distinct (for example, in $\C{R}0$, it is implicitly assumed that $i \ne j \ne k \ne i$).
For the special case $J = \emptyset$, the algorithm reduces to the standard formulation of the FCI algorithm \cite{Zhang2008_AI}.\footnote{Compared to the standard formulation of \cite{Zhang2008_AI}, which assumes no input nodes, we have adapted the skeleton search phase (by starting with a mixed graph that
contains no edges between input nodes, limiting the paths in the calculation of the
$\PosSEP_{H_0}(i,j)$ sets to output nodes only, and by including all input nodes, 
except $j$ itself, into the separating set). Furthermore, we added 
step~\ref{alg:fci:orientJ} to orient the edges between input and output nodes. 
The formulation of orientation rules $\C{R}0$, $\C{R}1$, $\C{R}3$, $\C{R}7$ is slightly different for this extended version
(as these would not be valid in case both $i, k \in J$), and the rest of the algorithm is unchanged.}
%We have omitted orientation rules $\C{R}5$--$\C{R}7$ of \cite{Zhang2008_AI}, since these are only necessary if selection bias is not ruled out a priori. 
%Those rules can yield edges of the form $\tuo, \out$ and $\tut$ that are not valid in $\sigma$-PAGs, and require the more general formalism of CPAGs (which is yet to be developed for the general case with possible cycles and input nodes).



\begin{algorithm}
\begin{algorithmic}[1]
  \Input Input node set $J$; output node set $V$; independence model $I$ over $V \given J$
  \Output mixed graph $H$ with input nodes $J$ and output nodes $V$
  \State $\lt H,\sepset \rt \leftarrow \mathtt{FCIskeleton}(J,V,I)$\label{alg:fci:skeleton}
  \ForEach{edge $j \ouo v$ in $H$ with $j \in J$, $v \in V$}\label{alg:fci:orientJ}\label{alg:fci:start_orient}
    \State orient $j \tuo v$
  \EndFor
%  \For{each edge $j \ouo v$ in $H$ with $j \in J$ and $v \in V$} \label{alg:fci:orientJ} \Comment{$\C{R}bk$}
%  \State orient it as $j \tuh v$ \Joris{Or $j \ouh v$?} 
%  \EndFor
%  \For{each unshielded triple $\lt i,j,k \rt$ in $H$ with $k \in V$} \Comment{$\C{R}0$}
%    \If{$j \notin \sepset(k,i)$}
%      \State orient it as $i \suh j \hus k$
%    \EndIf
%  \EndFor
  \Repeat \label{alg:fci:R0}
    \begin{description}
      \item[$\C{R}0$] if $i \sus j \sus k$ in $H$ with $i \notin J$, and $i$ and $k$ are not adjacent in $H$, then orient $i \suh j \hus k$ if $j \notin \sepset(\{i,k\})$ % \sepset(k,i) if ordered!
    \end{description}
  \Until{this orientation rule is not applicable}
  \Repeat
    \begin{description}
      \item[$\C{R}1$] if $i \suo j \hus k$ in $H$ with $i \notin J$, and $i$ and $k$ are not adjacent in $H$, then orient $i \hut j$
      \item[$\C{R}2$] if $i \tuh j \suh k$ or $i \suh j \tuh k$ in $H$, and $i \suo k$ in $H$, then orient $i \suh k$
      \item[$\C{R}3$] if $i \suh j \hus k$ and $i \suo l \ous k$ and $l \suo j$ in $H$ with $i \notin J$, and $i$ and $k$ are not adjacent in $H$, then orient $l \suh j$
      \item[$\C{R}4$] if $\lt i, j, q_1, \dots, q_n, k \rt$ is a discriminating path in $H$ for $j$, and if $i \suo j$ in $H$, then orient $i \hut j$ if $j \in \sepset(\{i,k\})$ and orient $i \huh j \huh q_1$ if $j \notin \sepset(\{i,k\})$ % \sepset(k,i) if ordered!
    \end{description}
  \Until{none of these orientation rules is applicable}
  \Repeat
    \begin{description}
      \item[$\C{R}5$] if $i \ouo j$ in $H$, and there is an uncovered circle path $i \ouo k \ouo \cdots \ouo l \ouo j$ in $H$ such that $i$ is not adjacent to $l$ and $j$ is not adjacent to $k$, then orient $i \tut k \tut \cdots \tut l \tut j \tut i$
    \end{description}
  \Until{this orientation rule is not applicable}
  \Repeat
    \begin{description}
      \item[$\C{R}6$] if $i \tut j \ous k$ in $H$, then orient $j \tus k$
      \item[$\C{R}7$] if $i \suo j \out k$ in $H$ with $i \notin J$, and $i$ and $k$ are not adjacent in $H$, then orient $i \sut j$
    \end{description}
  \Until{none of these orientation rules is applicable}
  \Repeat
    \begin{description}
      \item[$\C{R}8$] if $i \tuh j \tuh k$ in $H$, and $i \ouh k$ in $H$, then orient $i \tuh k$
      \item[$\C{R}9$] if $i \ouh k$, and $\pi = \lt i, j, \dots, k \rt$ is an uncovered possibly directed path in $H$ from $i$ to $k$ such that $j$ and $k$ are not adjacent in $H$, then orient $i \tuh k$
      \item[$\C{R}10$] if $i \ouh k$ in $H$, $j \tuh k \hut l$ in $H$, $\pi_1$ is a uncovered possibly directed path in $H$ from $i$ to $j$, and $\pi_2$ is a uncovered possibly directed path in $H$ from $i$ to $l$, then let $u_1$ be the node adjacent to $i$ on $\pi_1$ (possibly $u_1 = j$) and $u_2$ the node adjacent to $i$ on $\pi_2$ (possible $u_2 = l$); if $u_1 \ne u_2$, and $u_1$ and $u_2$ are not adjacent in $H$, then orient $i \tuh k$
    \end{description}
  \Until{none of these orientation rules is applicable}\label{alg:fci:end_orient}
  \end{algorithmic}
  \caption{Extended FCI Algorithm.\label{alg:fci}}
\end{algorithm}

\begin{Thm}[Extended FCI Soundness]\label{thm:fci_sound}
  The Extended FCI algorithm (Algorithm~\ref{alg:fci}) is \emph{sound}: if its input consists of the $\sigma$-independence model $\IM_\sigma(G \given S)$ of a CDMG $G$ given $S$, then its output will be a valid $\sigma$-PAG $H$ that represents $G$ given $S$.
\end{Thm}
\begin{proof}
  Let $G$ be a CDMG with input nodes $J$ and output nodes $V^+ = V \cup S$ and $I = \IM_\sigma(G \given S)$ its $\sigma$-independence model given $S$.

  The skeleton phase in line \ref{alg:fci:skeleton}, which invokes Algorithm~\ref{alg:fci_skeleton}, is sound (Theorem~\ref{thm:fci_skeleton_sound}).
  That is, it computes $H = \skel(\sigmaPAG(G \given S))$, and $\sepset$ will contain a separating set of $i$ and $j$ for every edge $i \ouo j$ absent in $H$ with $i \in V, j \in J \cup V, i \ne j$.
  The next step, line \ref{alg:fci:orientJ}, orients the edges between input and output nodes in $H$. 
  The result is then a valid $\sigma$-PAG that represents $G$ given $S$.
  The rest of the proof proceeds by induction.

  Now that the skeleton has been determined, the orientations (edge marks) of the edges will be deduced.
  We show for each of the orientation rules that under the assumption that the current $H$ is a valid $\sigma$-PAG that represents $G$ given $S$, applying the rule yields an updated $H$ that is still a valid $\sigma$-PAG that represents $G$ given $S$.
  We first exploit the modeling assumptions on the input nodes in line \ref{alg:fci:orientJ} to partially orient edges connecting input and output nodes.
  The soundness of this orientation step stems from Lemma~\ref{lem:inducing_path_out_of_input}.

  Some of the rules ($\C{R}1$, $\C{R}3$, $\C{R}5$, $\C{R}6$, $\C{R}7$, $\C{R}9$, $\C{R}10$) assume that rule $\C{R}0$ has been exhaustively applied, which is the reason that rule $\C{R}0$ is performed before the other orientation rules are performed.
  Additionally, rule $\C{R}6$ assumes that rule $\C{R}5$ has been exhaustively applied.

  In the following, we will always assume that the antecedent of the rule holds for a mixed graph $H$ that is a valid $\sigma$-PAG that represents CDMG $G$ given $S$.
  This implies, in particular, that if $v \hus w$ in $H$ or $v \ous w$ in $H$, then $v \notin J$.
  \begin{description}
    \item[$\C{R}0$] ``If $i \sus j \sus k$ in $H$, with $i \notin J$, and $i$ and $k$ are not adjacent in $H$, then orient $i \suh j \hus k$ if $j \notin \sepset(\{i,k\})$.''\\ % \sepset(k,i) if ordered!
      It follows from Proposition~\ref{prp:orienting_unshielded_triples} that $H$ still represents $G$ given $S$ after the orientation of the unshielded collider.
      %Note that there is no need to assume $j \in V$, because if $j \in J$, then it will be contained in $\sepset(\{i,k\})$.
    \item[$\C{R}1$] ``If $i \suo j \hus k$ in $H$ with $i \notin J$, and $i$ and $k$ are not adjacent in $H$, then orient $i \hut j$.''\\
%      \Joris{Note that $j,k \notin J$, but $i$ could be in $J$. The proof works in both cases and there is no risk of overwriting background knowledge.}
      We first show that $j \in \Anc_G(i)$.
      Since the triple $\lt i, j, k \rt$ is an unshielded triple in $H$ with $i \notin J$, but has not been oriented as a collider by $\C{R}0$, we conclude that $j \in \sepset(\{i,k\})$. % \sepset(k,i) if ordered!
      By Proposition~\ref{prp:orienting_unshielded_triples}, $j \in \Anc_G(\{i,k\} \cup S)$.
      Since $H$ represents $G$ and $j \hus k$ in $H$, $j \not\in \Anc_G(\{k\}\cup S)$. 
      Therefore, $j \in \Anc_G(i)$.
      Thus if we orient $i \sut j$, the mixed graph still represents $G$ given $S$.
      In particular, the orientation $i \tut j \hus k$ with $i, k$ non-adjacent cannot occur.
      Therefore, if we orient $i \hut j$, the resulting mixed graph $H$ will still represent $G$.
%      After orienting $j \tuh k$ in $H$, the resulting mixed graph $H$ still represents $G$.
    \item[$\C{R}2$] ``If $i \tuh j \suh k$ or $i \suh j \tuh k$ in $H$, and $i \suo k$ in $H$, then orient $i \suh k$.''\\
%      \Joris{Note that $j,k \notin J$, but $i$ might be in $J$. There is no risk of overwriting background knowledge.}
      By the antecedent of the rule, and since $H$ represents $G$, we have $k \notin \Anc_G(\{j\} \cup S)$.
      In case $i \tuh j \suh k$, we have $i \in \Anc_G(\{j\} \cup S)$, so if $k$ were in $\Anc_G(i)$, it would follow that $k \in \Anc_G(\{j\}\cup S)$, a contradiction.
      In case $i \suh j \tuh k$, we have $j \in \Anc_G(\{k\} \cup S)$, so if $k$ were in $\Anc_G(i)$, it would follow that $j \in \Anc_G(\{i\}\cup S)$, a contradiction.
      Hence, in both cases, we must have $k \not\in \Anc_G(i)$.
      In both cases, we also have $k \notin \Anc_G(S)$ because of the arrowhead on $j \suh k$.
      Therefore, after orienting $i \suh k$ in $H$, the resulting mixed graph still represents $G$.
    \item[$\C{R}3$] ``If $i \suh j \hus k$ and $i \suo l \ous k$ and $l \suo j$ in $H$ with $i \notin J$, and $i$ and $k$ are not adjacent in $H$, then orient $l \suh j$.''\\
%      \Joris{Note $j \notin J$, and $l \notin J$, and therefore also $i,k\notin J$. So there is no risk of overwriting background knowledge.}
      Since $\lt i, l, k \rt$ is an unshielded triple in $H$ with $i \notin J$ that was not oriented as a collider by $\C{R}0$, we must have that $l \in \Anc_G(\{i,k\} \cup S)$ by Proposition~\ref{prp:orienting_unshielded_triples}.
      Assume, for the sake of contradiction, that $j \in \Anc_G(\{l\} \cup S)$. 
      Then $j \in \Anc_G(\{i,k\} \cup S)$.
      This contradicts that $j \notin \Anc_G(\{i,k\} \cup S)$ from $i \suh j \hus k$ in $H$.
      Hence, $j \notin \Anc_G(l \cup S)$.
%      Furthermore, $j \notin \Anc_G(S)$ because of the arrowheads at $i \suh j \hus k$.
      After orienting $l \suh j$ in $H$, the resulting mixed graph still represents $G$.
    \item[$\C{R}4$] ``If $\lt i, j, q_1, \dots, q_n, k \rt$ is a discriminating path in $H$ for $j$, and if $i \suo j$ in $H$, then orient $i \hut j$ if $j \in \sepset(\{i,k\})$ and orient $i \huh j \huh q_1$ if $j \notin \sepset(\{i,k\})$.''\\ % \sepset(k,i) if ordered!
%      \Joris{This already seems OK: only $j, k$ could be in $J$, and for sure $i \in V$.}
%      \Joris{Note that $q_1,\dots,q_n \notin J$, $k \notin J$, but $i$ and $j$ might be in $J$. There is the risk of overwriting background knowledge, so we should change the rule (similar to $\C{R}0$). However, since $J\sm\{k,i\}$ is part of $\sepset(k,i)$, there is actually no risk of overwriting background knowledge.}
      It follows immediately from Proposition~\ref{prp:orienting_discriminating_paths} that applying this rule yields an updated mixed graph that still represents $G$.
    \item[$\C{R}5$] ``If $i \ouo j$ in $H$, and there is an uncovered circle path $i \ouo k \ouo \cdots \ouo l \ouo j$ in $H$ such that $i$ is not adjacent to $l$ and $j$ is not adjacent to $k$, then orient $i \tut k \tut \cdots \tut l \tut j \tut i$.''\\
      There is an uncovered cycle consisting of (at least 4) $\ouo$ edges in $H$. 
      Lemma~\ref{lem:detecting_selection_bias} implies that each node on the uncovered cycle must be in $\Anc_G(S)$.
      Hence we can orient all edges on the cycle as undirected, and this yields a mixed graph that still represents $G$ given $S$.
    \item[$\C{R}6$] ``If $i \tut j \ous k$ in $H$, then orient $j \tus k$.''\\
%      If $i \in \Sc_G(j)$ then orienting instead $j \hus k$ would give a contradiction, as then we also need $i$ and $k$ adjacent by the little lemma. No: $i$ and $k$ can be adjacent!
%      If $i \notin \Sc_G(j)$ then (acyclic case!) orienting instead $j \hus k$ would give a contradiction.
      Only $\C{R}5$ could have introduced the undirected edge.
      In that case, we know that both $i$ and $j$ are in $\Anc_G(S)$ from Lemma~\ref{lem:detecting_selection_bias}. 
      Hence we can orient $j \tus k$ and the updated mixed graph will still represent $G$ given $S$.
    \item[$\C{R}7$] ``If $i \suo j \out k$ in $H$ with $i \notin J$, and $i$ and $k$ are not adjacent in $H$, then orient $i \sut j$.''\\
      Suppose after orienting $i \suh j \out k$, the mixed graph would still represent $G$ given $S$.
      If $j \in \Anc_G(\{k\} \cup S)$, then we could further orient $i \suh j \tut k$ and the mixed graph would still represent $G$ given $S$, yielding a contradiction because
      an unshielded triple of that form cannot occur. 
      So $j \notin \Anc_G(\{k\} \cup S)$, and we can orient $i \suh j \hut k$ to obtain a mixed graph that still represents $G$ given $S$.
      But then we have an unshielded collider $(i,j,k)$ with $i \notin J$ that should have been oriented by $\C{R}0$, another contradiction.
      Hence we can orient $i \sut j$ and the resulting mixed graph will still represent $G$ given $S$.
    \item[$\C{R}8$] ``If $i \tuh j \tuh k$ in $H$, and $i \ouh k$ in $H$, then orient $i \tuh k$.''\\
%      \Joris{$j,k\notin J$, $i$ might be. No risk in overwriting background knowledge.}
      %Clearly, $i \in \Anc_G(\{j\} \cup S)$ and $j \in \Anc_G(\{k\} \cup S)$ implies $i \in \Anc_G(\{k\} \cup S)$.
      It follows immediately from Lemma~\ref{lem:mixed_graph_represents_CDMG} that $i \in \Anc_G(\{k\} \cup S)$.
      Thus, applying this rule yields an updated mixed graph that still represents $G$.
    \item[$\C{R}9$] ``If $i \ouh k$, and $\pi = \lt i, j, \dots, k \rt$ is an uncovered possibly directed path in $H$ from $i$ to $k$ such that $j$ and $k$ are not adjacent in $H$, then orient $i \tuh k$.''\\
%      \Joris{$j, \dots, k \notin J$, $i \notin J$. No risk to overwrite background knowledge.}
      %Suppose $i \notin \Anc_G(\{k\} \cup S)$.
      Note that $j \notin J$ as either $i \suh j$ or $i \suo j$.
      The unshielded triple $\lt j, i, k \rt$, with $j \notin J$, was not oriented as a collider by $\C{R}0$, and therefore $i \in \Anc_G(\{j,k\} \cup S)$.
      If $i \in \Anc_G(j)$, then Lemma~\ref{lem:upd_path} then states that $i \in \Anc_G(k)$.
      Therefore, we conclude that $i \in \Anc_G(\{k\} \cup S)$.
      Applying the rule therefore yields an updated mixed graph that still represents $G$.
    \item[$\C{R}10$] ``Suppose that $i \ouh k$ in $H$, $j \tuh k \hut l$ in $H$, $\pi_1$ is a uncovered possibly directed path in $H$ from $i$ to $j$, and $\pi_2$ is an uncovered possibly directed path in $H$ from $i$ to $l$. Let $u_1$ be the node adjacent to $i$ on $\pi_1$ (possibly $u_1 = j$) and $u_2$ the node adjacent to $i$ on $\pi_2$ (possibly $u_2 = l$). If $u_1 \ne u_2$, and $u_1$ and $u_2$ are not adjacent in $H$, then orient $i \tuh k$.''\\
%      \Joris{$i \notin J$, $k \notin J$, $j \notin J$, $l \notin J$, $u_1 \notin J$, $u_2 \notin J$. No risk to overwrite background knowledge.}
      %We will assume that $i \notin \Anc_G(\{k\} \cup S)$ and derive a contradiction.
      Because $u_1$ lies on a possibly directed path starting at $i$, it cannot be in $J$; the same holds for $u_2$.
      The unshielded triple $\lt u_1, i, u_2 \rt$, with $u_1, u_2 \notin J$, was not oriented by rule $\C{R}0$, which implies that $i \in \Anc_G(\{u_1,u_2\} \cup S)$. 
      %By assumption, $i \notin \Anc_G(S)$, so $i \in \Anc_G(u_1)$ or $i \in \Anc_G(u_2)$.
      If $i \in \Anc_G(u_1)$, Lemma~\ref{lem:upd_path} gives that $i \in \Anc_G(j)$.
      If $i \in \Anc_G(u_2)$, Lemma~\ref{lem:upd_path} gives that $i \in \Anc_G(l)$.
%      If $u_1 = j$, then we conclude that $i \in \Anc_G(k)$, a contradiction.
%      So assume $u_1 \ne j$.
%      We distinguish two cases: $u_1 \in \Anc_G(i)$ and $u_1 \notin \Anc_G(i)$.
%
%      If $u_1 \in \Anc_G(i)$, then $u_1 \in \Sc_G(i)$, and then $k \huh u_1$ must be in $H$ by Lemma~\ref{lem:induced_path_into_scc}.
%      Denote the node on $\pi_1$ next to $u_1$ by $v_1$ (possibly $v_1 = j$).
%      As $v_1$ is not adjacent to $i$, by Lemma~\ref{lem:induced_path_into_scc}, $u_1 \notin \Anc_G(\{v_1\} \cup S)$ would give a contradiction, and hence $u_1 \in \Anc_G(\{v_1\} \cup S)$. 
%      Because of the arrowhead in $k \huh u_1$ at $u_1$, $u_1 \notin \Anc_G(S)$.
%      Hence $u_1 \in \Anc_G(v_1)$, and we can orient $u_1 \tus v_1$.
%      If $v_1 \in \Anc_G(\{u_1\} \cup S)$, then $v_1 \in \Sc_G(u_1)$ (as $v_1 \in \Anc_G(S)$ would contradict $u_1 \notin \Anc_G(S)$).
%      But then $i$ and $v_1$ should be adjacent \Joris{lemma?}.
%      Hence we can orient $u_1 \tuh v_1$.
%      By Lemma~\ref{lem:upd_path}, we can now orient all remaining edges on $\pi_1$ as directed towards $j$, and conclude that $u_1 \in \Anc_G(j)$, and therefore also $i \in \Anc_G(k)$.
%
%      If $u_1 \notin \Anc_G(i)$, then also $u_1 \notin \Anc_G(S)$ (otherwise $i \in \Anc_G(u_1) \subseteq \Anc_G(S)$), and we can orient $i \tuh u_1$.
%      By Lemma~\ref{lem:upd_path}, we can then orient all remaining edges on $\pi_1$ as directed towards $j$, and conclude that $u_1 \in \Anc_G(j)$, and therefore also $i \in \Anc_G(k)$.
      In both cases, $i \in \Anc_G(\{k\} \cup S)$, because $j \tuh k \hut l$ in $H$.
      So we conclude that $i \in \Anc_G(\{k\} \cup S)$ and can orient $i \tuh k$ to obtain an updated mixed graph that still represents $G$ given $S$.
  \end{description}
  Since each of these orientation rules leaves the skeleton (adjacencies) of $H$ invariant,
  %and each orientation of an edge mark leads to a valid edge,
  and after each orientation rule, $H$ still represents $G$,
  $H$ remains a valid $\sigma$-PAG throughout the orientation phase
  by Proposition~\ref{prp:sigmapag_represents_cdmg_is_valid}. 
\end{proof}
%We have here omitted orientation rules $\C{R}5$--$\C{R}7$ since these are only necessary if selection bias is not ruled out a priori. 
%Those rules can yield edges of the form $\tuo, \out$ and $\tut$ that are not valid in $\sigma$-PAGs, and require the more general formalism of PAGs extended with input nodes (which is yet to be written down).

The following two lemmata are applicable after the initial phase of the extended FCI algorithm, once rule $\C{R}0$ has been exhaustively applied.
The first concerns uncovered circle paths, the second uncovered possibly directed paths.
\begin{Lem}\label{lem:detecting_selection_bias}
  Let $H$ be a mixed graph that represents $G$ given $S$ in which rule $\C{R}0$ has been exhaustively applied.
  If $i \ouo j$ in $H$, and there is an uncovered circle path $i \ouo k \ouo \cdots \ouo l \ouo j$ in $H$ such that $i$ is not adjacent to $l$ and $j$ is not adjacent to $k$, then every node on the path is in $\Anc_G(S)$.
\end{Lem}
\begin{proof}
  There is an uncovered cycle consisting of (at least 4) $\ouo$ edges in $H$. 
  Note that none of the nodes on the cycle can be in $J$, as each node has two or more circle edge marks.
  Suppose we orient one of the circles into an arrowhead and the new mixed graph still represents $G$ given $S$.
  Then we could make use of rule $\C{R}1$ repeatedly to orient the whole cycle as a directed cycle, and the resulting mixed graph should still represent $G$ given $S$.
  However, that would be a contradiction, since it contains a directed cycle.
  Hence, any circle edge mark on the cycle that we orient as an arrowhead would yield a mixed graph that no longer represents $G$ given $S$.
  In other words, each node on the cycle must be ancestor in $G$ of its neighboring node, or of the selection set $S$.
  This implies that every node on the cycle must be in $\Anc_G(S)$.
  Indeed, if a given node on the cycle is not in $\Anc_G(S)$, then it must be ancestor in $G$ of its neighbors.
  Its neighbors then cannot be in $\Anc_G(S)$ either, and therefore must be ancestor in $G$ of their neighbors.
  But then we would have an unshielded triple of nodes that are in the same strongly connected component of $G$: a contradiction with Lemma~\ref{lem:sigma-inducing_path_in_scc}.
\end{proof}

\begin{Lem}\label{lem:upd_path}
Let $H$ be a mixed graph that represents $G$ given $S$ in which rule $\C{R}0$ has been exhaustively applied.
Let $v_1 \sus \dots \sus v_n$ be an uncovered possibly directed path from $v_1$ to $v_n$ (with $n \ge 2$) in $H$.
%If $v_1 \in \Anc_G(v_2 \cup S)$ and $v_2 \notin \Anc_G(v_1 \cup S)$, 
If there is an edge $k \suh v_1$ in $H$ and if $v_1 \in \Anc_G(\{v_2\} \cup S)$, then we conclude that $v_1 \in \Anc_G(v_n)$.
\end{Lem}
\begin{proof}
  Note that $v_1, v_2, \dots, v_n$ are all not in $J$, because each of them must have an arrow head or circle edge mark.
  We can orient $v_1 \tus v_2$ because of the assumption that $v_1 \in \Anc_G(\{v_2\} \cup S)$.
  Note that $v_1 \in \Anc_G(v_2)$ because $v_1 \notin \Anc_G(S)$ due to the arrowhead at $v_1$ on the edge $k \suh v_1$.
  If $n = 2$, then we immediately conclude that $v_1 \in \Anc_G(v_n)$.
  So assume $n > 2$.
  We distinguish two cases: $v_2 \in \Anc_G(v_1)$ and $v_2 \notin \Anc_G(v_1)$.

  If $v_2 \notin \Anc_G(v_1)$, we can orient $v_1 \tuh v_2$, since $v_2 \in \Anc_G(S)$ would imply $v_1 \in \Anc_G(S)$, a contradiction.
  By the same reasoning as in the proof of rule $\C{R}1$, we conclude that $v_2 \in \Anc_G(v_3)$, and that we can orient $v_2 \tuh v_3$.
  We can iterate this reasoning subsequently on all remaining edges on the uncovered possibly directed path to deduce that $v_i \in \Anc_G(v_{i+1})$ and we can orient $v_i \tuh v_{i+1}$, for $i = 3, \dots, n-1$.
  This leads to the conclusion that $v_1 \in \Anc_G(v_n)$.

  If $v_2 \in \Anc_G(v_1)$, that is $v_2 \in \Sc_G(v_1)$, then we can orient $v_1 \tut v_2$. 
  By Lemma~\ref{lem:inducing_path_into_scc}, this implies that $k$ and $v_2$ must be adjacent in $H$ as well. 
  This edge can be oriented as $k \suh v_2$, because $v_2 \in \Anc_G(\{k\} \cup S)$ would imply $v_1 \in \Anc_G(\{k\} \cup S)$, a contradiction.
  The assumption $v_2 \notin \Anc_G(\{v_3\} \cup S)$ gives a contradiction: we could then orient $v_2 \hus v_3$, obtaining an unshielded triple $v_1 \tut v_2 \hus v_3$.
  Hence $v_2 \in \Anc_G(\{v_3\} \cup S)$, and we can orient $v_2 \tus v_3$. 
  Because of the arrowhead in $k \suh v_2$ at $v_2$, $v_2 \notin \Anc_G(S)$, and hence $v_2 \in \Anc_G(v_3)$.
  If $v_3 \in \Anc_G(\{v_2\} \cup S)$, then $v_3 \in \Sc_G(v_2)$ (as $v_3 \in \Anc_G(S) \sm \Anc_G(v_2)$ would contradict $v_2 \notin \Anc_G(S)$).
  But then $v_1$ and $v_3$ must lie in the same strongly connected component of $G$, and should therefore be adjacent in $H$ by Lemma~\ref{lem:sigma-inducing_path_in_scc}, which they are not.
  Hence $v_3 \notin \Anc_G(\{v_2\} \cup S)$ and we can orient $v_2 \tuh v_3$.
  The reasoning now proceeds as in the previous case, and leads to the conclusion that $v_2 \in \Anc_G(v_n)$.
  Because $v_1 \in \Anc_G(v_2)$, also $v_1 \in \Anc_G(v_n)$.
\end{proof}

The soundness of the Extended FCI algorithm immediately implies its consistency when using consistent conditional independence tests.
\begin{Cor}[Extended FCI Consistency]\label{cor:fci_consistent}
%  The Extended FCI algorithm (Algorithm~\ref{alg:fci}) is \emph{consistent}: if its input consists of the $\sigma$-independence model $\IM_\sigma(G \given S)$ of a CDMG $G$ given $S$, then its output will be a valid $\sigma$-PAG $H$ that represents $G$ given $S$.
  Let $M = \lt \Sin, \Send^+, \Srnd, \CX, P, f \rt$ be a simple SCM with endogenous variables $V^+ = V \dcup S \dcup L$ (note that we allow additional latent endogenous variables $L$ here).
  Let $\xi_S \subseteq \CX_S$ be a measurable set with $\Prb_M(X_S \in \xi_S \given \doit(X_J = x_J)) > 0$ for all $x_j \in \CX_J$.
  Assume that we have access to infinitely many samples distributed according to the marginal Markov kernel of $M$ on $V$ after selecting on $S$:
    $$P_M(X_V \mid X_S \in \xi_S, \doit(X_J=x_J)).$$
  Assume also that the following faithfulness assumption holds:
    $$\IM(P_{M}(X_V \mid X_S \in \xi_S, \doit(X_J=x_J))) = \IM_\sigma(G_{V \cup S \given J}(M) \given S).$$
  When using asymptotically consistent conditional independence tests (on the i.i.d.\ samples of the Markov kernel $P_M(X_V \mid X_S \in \xi_S, \doit(X_J=x_J))$), 
  the Extended FCI algorithm (Algorithm~\ref{alg:fci}) provides an asymptotically consistent estimate $\hat{H}$ of the $\sigma$-PAG $\FCI(\IM_\sigma(G_{V \cup S \given J}(M) \given S))$, which represents $G_{V \cup S \given J}(M)$ given $S$.
\end{Cor}
\begin{proof}
  The asymptotic consistency of the conditional independence tests means that the probability of a wrong test result (either Type I or Type II error) vanishes asymptotically.
  Since $\IM(P_{M}(X_V \mid X_S \in \xi_S, \doit(X_J=x_J)))$ consists of finitely many conditional independence statements, the test results will agree completely with this conditional independence model with arbitrarily high probability given sufficiently many samples.
  By the faithfulness assumption, this conditional independence model agrees with $\IM_\sigma(G_{V \cup S \given J}(M) \given S)$.
  By the soundness of the Extended FCI algorithm (Theorem~\ref{thm:fci_sound}), the $\sigma$-PAG that FCI outputs will equal $\FCI(\IM_\sigma(G_{V \cup S \given J}(M) \given S))$ with arbitrarily high probability given sufficiently many samples.
\end{proof}
Because conditional independence tests are not uniformly consistent (as there is no upper bound on the number of samples needed to distinguish an arbitrarily weak dependence from an independence, without additional assumptions), also the Extended FCI algorithm is not uniformly consistent. 
In other words, it is not known in advance how many samples will be needed to yield a reliable result.


\subsection{Completeness}

For acyclic $G$ without input nodes (that is, if $G$ is an ADMG), the FCI algorithm was shown to be complete \cite{Zhang2008_AI} in the sense that all edge marks that could possibly be oriented based on the information in $\IM_\sigma(G \given S)$ will be oriented.
Using results on the characterization of Markov equivalence classes of maximal ancestral graphs \cite{Ali++2005}, it can additionally be shown that the $\sigma$-PAG output by FCI represents the Markov equivalence class of $G$ with respect to $V$ in case $G$ is an ADMG.
By employing acyclifications, these results have been extended to cyclic $G$ \cite{MooijClaassen_UAI_20}, still without input nodes, and with the additional assumption of no selection bias ($S = \emptyset$).
The known completeness results have very long proofs, and we will therefore not provide these here.
Instead, we will only formulate these results, and refer the interested reader to the original papers for the proofs.


We make the following assumption in order to state the known completeness results.\footnote{While we believe that it is straightforward to extend the completeness results to allow for both cycles and selection bias, it is known that the Extended FCI algorithm (Algorithm~\ref{alg:fci}) is incomplete when also allowing for input nodes.}
\begin{Ass}\label{ass:fci_completeness_assumption}
  Given node set $V$, let:
  \begin{itemize}
    \item $\C{G}_V$ be the set of all pairs $(G,S)$ of acyclic DMGs $G$ with output nodes $V^+ = V \dcup S$ for some disjoint set $S$ (`acyclicity'), or
    \item $\C{G}_V$ be the set of all pairs $(G,\emptyset)$ of DMGs $G$ with output nodes $V^+ = V \dcup \emptyset$ (`no selection bias').
  \end{itemize}
  In both cases, all DMGs in $\C{G}_V$ have $J = \emptyset$ (`no input nodes').
\end{Ass}

In the absence of input nodes (for $J = \emptyset$), we can interpret FCI as a mapping that maps an independence model (on $V$) to a mixed graph (with nodes $V$).
By Theorem~\ref{thm:fci_sound}, it maps $\IM_\sigma(G\given S)$, the $\sigma$-independence model of a DMG $G$ given $S$, to a $\sigma$-PAG $\FCI(\IM_\sigma(G \given S))$ that represents $G$ given $S$.
Additionally, the following (incomplete) completeness results are known.
\begin{Thm}[Some FCI completeness results]\label{thm:fci_complete}
Under Assumption~\ref{ass:fci_completeness_assumption}, the Extended FCI algorithm (Algorithm~\ref{alg:fci}) is:
  \begin{enumerate}[label=(\roman*)]
    \item \emph{arrowhead complete}: for all $(G,S) \in \C{G}_V$, for all $i \ne j \in V$: there is an arrowhead $i \hus j$ in $\FCI(\IM_\sigma(G \given S))$ if $i \notin \Anc_{\tilde{G}}(\{j\} \cup S)$ for all $(\tilde{G},\tilde{S}) \in \C{G}_V$ with $\IM_\sigma(\tilde{G} \given \tilde{S}) = \IM_\sigma(G \given S)$;
    \item \emph{tail complete}: for all $(G,S) \in \C{G}_V$, for all $i \ne j \in V$: there is a tail $i \tus j$ in $\FCI(\IM_\sigma(G \given S))$ if $i \in \Anc_{\tilde{G}}(\{j\} \cup S)$ for all $(\tilde{G},\tilde{S}) \in \C{G}_V$ with $\IM_\sigma(\tilde{G} \given \tilde{S}) = \IM_\sigma(G \given S)$;
    \item \emph{Markov complete}: for all $(G_1,S_1) \in \C{G}_V$ and $(G_2,S_2) \in \C{G}_V$: $\IM_\sigma(G_1 \given S_1) = \IM_\sigma(G_2 \given S_2)$ if and only if $\FCI(\IM_\sigma(G_1 \given S_1)) = \FCI(\IM_\sigma(G_2 \given S_2))$.
  \end{enumerate}
\end{Thm}
\begin{proof}
  The first two claims are proved in \cite{Zhang2008_AI} under the additional assumption of acyclicity.
  In \cite{MooijClaassen_UAI_20} it is explained how the characterization of the Markov equivalence classes of \cite{Ali++2005} can be used to then prove the third claim under that additional assumption.
  Furthermore, in \cite{MooijClaassen_UAI_20} it is shown how to generalize these results to the cyclic setting by employing acyclifications, but only under the additional assumption of no selection bias.
\end{proof}

Arrowhead and tail completeness express that the $\sigma$-PAG output by FCI is maximally oriented: any arrowhead or tail that could possibly be deduced from $\IM_\sigma(G \given S)$, will have been oriented as such in the $\sigma$-PAG.
The soundness and Markov completeness properties together imply that the $\sigma$-PAG output by FCI, when given as input the $\sigma$-independence model of a directed mixed graph given some set of latent selection nodes, represents the $\sigma$-Markov equivalence class of $G$ with respect to the observed nodes.
In other words, FCI provides a \emph{graphical characterization} of the $\sigma$-Markov equivalence class.
